% ============================================================
%  第二章：偏好、效用与显示偏好（精讲版）
%  依据：Suzumura《Choice, Preferences, and Procedures》Chapter 1 与课程材料
%       + Suzumura《Choice, Preferences, and Procedures》Chapter 1
%  编译：LuaLaTeX（两遍）
% ============================================================
% 第 2 章 统一书稿前言包装
\newcommand{\lectureno}{第二章}
\newcommand{\lectureshort}{偏好、效用与显示偏好}
% ============================================================
%  微观经济学理论课程讲义 —— 统一课程前言（四讲共用）
%  版式参考 Nolan Miller, Notes on Microeconomic Theory (SWP/notes3 风格)
%  编译：LuaLaTeX（两遍）
%  各讲 preamble.tex 只需设定 \lectureno / \lectureshort 后 \input 本文件。
% ============================================================
\documentclass[11pt,a4paper,fontset=none,oneside,openany]{ctexbook}

% ---------- 版式 ----------
\usepackage{geometry}
\geometry{
  a4paper,
  top=2.8cm, bottom=2.8cm,
  left=2.6cm, right=2.6cm,
  headsep=0.85cm, footskip=1.25cm
}
\setlength{\headheight}{14pt}

% ---------- 字体（全部矢量嵌入，保证 PDF 高清晰度）----------
% 注意：本机 macOS 已移除 Kaiti SC，CJK 斜体统一回退到 Songti SC；
%        Songti SC / Heiti SC 经 OSFONTDIR 指向 Supplemental 后由 luaotfload 解析。
\usepackage{fontspec}
\setmainfont{TeX Gyre Termes}[Ligatures=TeX]
\setsansfont{TeX Gyre Heros}[Ligatures=TeX,Scale=MatchLowercase]
\setmonofont{TeX Gyre Cursor}[Scale=MatchLowercase]
\setCJKmainfont{Songti SC}[BoldFont={Songti SC Bold},ItalicFont={Songti SC}]
\setCJKsansfont{Heiti SC}[BoldFont={Heiti SC}]

% ---------- 数学 ----------
\usepackage{amsmath}
\usepackage{mathtools}
\usepackage{amsthm}
\usepackage{unicode-math}
\setmathfont{TeX Gyre Termes Math}
\allowdisplaybreaks[2]
\setlength{\abovedisplayskip}{9pt plus 3pt minus 3pt}
\setlength{\belowdisplayskip}{9pt plus 3pt minus 3pt}

% ---------- 定理类环境（书稿格式：按章编号，如“定理 1.1”）----------
\theoremstyle{plain}
\newtheorem{theorem}{定理}[chapter]
\newtheorem{proposition}[theorem]{命题}
\newtheorem{lemma}[theorem]{引理}
\newtheorem{corollary}[theorem]{推论}
\theoremstyle{definition}
\newtheorem{definition}[theorem]{定义}
\newtheorem{condition}[theorem]{条件}
\newtheorem{example}[theorem]{例}
\theoremstyle{remark}
\newtheorem{remark}[theorem]{注}
% 公式、图、表按章编号（书稿格式：式 (1.1)、图 1.1、表 1.1）
\numberwithin{equation}{chapter}
\numberwithin{figure}{chapter}
\numberwithin{table}{chapter}

% 定理类环境标题后的标点：amsthm 默认加英文 "."，此处统一改为中文句号。
% 原理：amsthm 的 \@thm 每次开环境都执行 \thm@headpunct{.}，而样式宏在其之后展开，
% 故把 \thm@headpunct{。} 追加到三个样式宏末尾即可覆盖（用 \AtBeginDocument 以免被 ctex 覆盖）。
\makeatletter
\AtBeginDocument{%
  \let\WBth@plain\th@plain
  \let\WBth@definition\th@definition
  \let\WBth@remark\th@remark
  \def\th@plain{\WBth@plain\thm@headpunct{。}}%
  \def\th@definition{\WBth@definition\thm@headpunct{。}}%
  \def\th@remark{\WBth@remark\thm@headpunct{。}}%
}
\makeatother

% 无编号小标题（仿 notes3 的 “A Note on ...”）
\newcommand{\notehead}[1]{\par\medskip\noindent{\bfseries #1}\par\smallskip}

% ---------- 表格 ----------
\usepackage{booktabs}
\usepackage{array}
\renewcommand{\arraystretch}{1.25}
\setlength{\tabcolsep}{8pt}

% ---------- 图形 ----------
\usepackage{graphicx}
\usepackage{tikz}
\usepackage{pgfplots}
\pgfplotsset{compat=1.18}
\usetikzlibrary{arrows.meta,positioning,shapes.geometric,calc,patterns,decorations.pathreplacing,decorations.markings}
\usepackage{float}
\usepackage{caption}
\captionsetup{font=small,labelfont=bf,labelsep=quad,skip=4pt}

% 统一图形配色（黑白模式：黑/灰梯度，供所有 TikZ 图复用）
% —— 全局黑白化：原彩色调色板统一改为黑/灰，保证打印与投影均可读 ——
\definecolor{AxisBlue}{RGB}{0,0,0}        % 轴线 / 主曲线：黑
\definecolor{SoftBlue}{RGB}{130,130,130}  % 节点填充：中灰
\definecolor{LightRed}{RGB}{224,224,224}  % 浅填充：浅灰
\definecolor{DeepBlue}{RGB}{0,0,0}        % 描边：黑
\definecolor{MidBlue}{RGB}{70,70,70}      % 次曲线 / 网格：深灰
\tikzset{
  >=Stealth,
  every picture/.append style={font=\small},
  axisline/.style={thick,color=black,-Stealth},
  axislabel/.style={color=black},
  fillsoft/.style={fill=LightRed,fill opacity=0.6,draw=black,thick},
  nodeblue/.style={circle,fill=black!85,draw=black,thick,inner sep=2pt},
  nodewhite/.style={circle,fill=white,draw=black,thick,inner sep=2pt},
  arrblue/.style={->, thick, color=black},
  boxsoft/.style={draw=black, rounded corners=2pt, align=center, inner sep=5pt, fill=black!6},
  % 多曲线 / 多集合区分：以线型 + 标记代替颜色
  curveA/.style={thick,black,solid},
  curveB/.style={thick,black,dashed},
  curveC/.style={thick,black,dotted},
  markA/.style={mark=*,mark size=2.2pt,color=black},
  markB/.style={mark=square*,mark size=2.0pt,color=black},
  markC/.style={mark=triangle*,mark size=2.4pt,color=black},
  % 图内区域 / 点标签（白底避免与图形重叠）
  figlabel/.style={font=\scriptsize,fill=white,inner sep=1pt,outer sep=1pt,text=black},
}
% 图形说明（黑白模式下对图形区域 / 曲线 / 关键点的解释，置于图形下方）
\newcommand{\figexplan}[1]{%
  \par\smallskip\noindent{\small\itshape 图形说明：#1}\par\smallskip}

% ---------- 页眉页脚（书稿格式：左=书名，右=章次与章名）----------
\usepackage{fancyhdr}
% 章标题：居中、大号（书稿章首样式）
\CTEXsetup[format={\centering\LARGE\bfseries}]{chapter}
\CTEXsetup[beforeskip=16pt,afterskip=14pt]{chapter}
\CTEXsetup[format={\large\bfseries}]{section}
\CTEXsetup[format={\normalsize\bfseries}]{subsection}
\pagestyle{fancy}
\fancyhf{}
\providecommand{\booktitle}{微观经济学理论：可行集、偏好与效用}
\providecommand{\lectureno}{第一章}
\providecommand{\lectureshort}{课程讲义}
\fancyhead[L]{\small\itshape \booktitle}
\fancyhead[R]{\small \lectureno\quad \lectureshort}
\fancyfoot[C]{\small \thepage}
\renewcommand{\headrulewidth}{0.4pt}
\fancypagestyle{plain}{\fancyhf{}\fancyfoot[C]{\small \thepage}\renewcommand{\headrulewidth}{0pt}}

% ---------- 正文细节 ----------
\usepackage{setspace}
\setstretch{1.12}
\usepackage{microtype}
\tolerance=900
\emergencystretch=2em
\raggedbottom

% ---------- 超链接与 PDF 书签 ----------
\usepackage{hyperref}
\hypersetup{
  colorlinks=true,
  linkcolor=black,
  urlcolor=black,
  citecolor=black,
  pdfencoding=auto,
  bookmarks=true,
  pdftitle={\booktitle\quad \lectureno：\lectureshort},
  pdfsubject={微观经济学理论书稿}
}
% 章节标题中常见数学符号的安全书签映射（避免 "Improper alphabetic constant"）
\pdfstringdefDisableCommands{%
  \def\leftrightarrow{<->}%
  \def\star{*}%
  \def\S{S}%
  \def\succeq{\text{≽}}%
  \def\succ{\text{≻}}%
  \def\sim{\text{∼}}%
}

% ---------- 常用记号 ----------
\newcommand{\R}{\mathbb{R}}
\newcommand{\N}{\mathbb{N}}
\newcommand{\Z}{\mathbb{Z}}
\newcommand{\Q}{\mathbb{Q}}
\newcommand{\succeqq}{\succeq}
\newcommand{\argmax}{\operatorname*{arg\,max}}
\newcommand{\argmin}{\operatorname*{arg\,min}}
\newcommand{\st}{\text{s.t.}\ }

\setcounter{chapter}{1}


\begin{document}

\chapter{偏好、效用与显示偏好：理性选择理论的形式化基础}

\noindent\textbf{本章任务。}第一章我们建立了选择理论的问题意识：什么是理性、如何用模型表达它。
本章进入形式化核心，沿着\textbf{三条线索}展开并把它们缝合为一个整体：

\begin{enumerate}
  \item \textbf{偏好路径}：偏好关系 $\succeq$ 需要满足哪些公理才算“理性”？（完备性、传递性及其弱化）
  \item \textbf{效用路径}：什么时候存在效用函数 $u(\cdot)$ 把偏好“数字化”？（有限集表示定理、词典序反例、序数性）
  \item \textbf{显示偏好路径}：如果只能\textbf{观察选择行为}，如何反推偏好、检验理性？（WARP/SARP/HARP，HARP $\Leftrightarrow$ 理性化）
\end{enumerate}

\noindent 三条线索的汇合点是本章的两个\textbf{枢纽定理}：
\textbf{效用表示定理}（偏好 $\to$ 数字）与 \textbf{HARP 定理}（选择行为 $\to$ 偏好，Richter--Hansson--Suzumura）。
最后以 Suzumura 的任意定义域理论（WCA/SCA 同余公理、G-理性与 M-理性）收束，说明这套框架在多大程度上可以不依赖预算集结构而成立。

\medskip
\noindent\textbf{主要参考}：Suzumura, K., \emph{Choice, Preferences, and Procedures: A Rational Choice Theoretic Approach}（Chapter 1: Rational Choice and Revealed Preference）；Levin (Stanford) \emph{Choice Theory} 讲义。

\tableofcontents
\clearpage

% ============================================================
\section{偏好关系的公理体系}
% ============================================================

\subsection{二元关系与偏好的分解}

\begin{definition}[二元关系、弱偏好]
设 $X$ 为备选方案的全集（universal set of alternatives）。$X$ 上的一个\textbf{二元关系}
$R \subseteq X \times X$。记全体二元关系的集合为 $\mathcal{B}$。我们把 $(x,y)\in R$ 读作
“$x$ 至少和 $y$ 一样好”（weakly preferred），此时称 $R$ 为 $X$ 上的\textbf{弱偏好关系}。
\end{definition}

一个弱偏好关系 $R$ 天然地拆成两块：

\begin{figure}[H]
\centering
\begin{tikzpicture}[font=\scriptsize,scale=1.15,transform shape]
  \draw[thick,fill=gray!15] (0,0) rectangle (2.6,2.9);
  \node[font=\small\bfseries] at (1.3,2.56) {(a) 关系 $R$ 的单元格};
  \foreach \x in {0.5,1.0,1.5,2.0} \draw[thin] (\x,0.62) -- (\x,2.12);
  \foreach \y in {0.62,1.12,1.62,2.12} \draw[thin] (0.5,\y) -- (2.0,\y);
  \foreach \c/\r in {0/0,0/1,1/0,1/2,2/1} \fill[gray!55] ({0.5+0.5*\c},{0.62+0.5*\r}) rectangle ({1.0+0.5*\c},{1.12+0.5*\r});
  \fill[gray!15] (1.5,0.62) rectangle (2.0,1.12);
  \node at (1.25,0.34) {列 $=$ 第一个元素，行 $=$ 第二个元素};
  \fill[gray!55] (0.30,-0.22) rectangle (0.52,0.00);
  \node[anchor=west,fill=white,inner sep=1pt] at (0.60,-0.11) {$I(R)$：双向成立};
  \fill[gray!15] (0.30,-0.62) rectangle (0.52,-0.40);
  \node[anchor=west,fill=white,inner sep=1pt] at (0.60,-0.51) {$P(R)$：仅正向成立};
  \begin{scope}[xshift=5.4cm]
    \draw[thick,fill=gray!15] (0,0) rectangle (3.2,2.9);
    \node[font=\small\bfseries] at (1.6,2.56) {(b) 在 $X=\{a,b,c\}$ 上的读法};
    \fill (0.6,1.62) circle (1.5pt) node[below=2pt] {$a$};
    \fill (1.6,1.62) circle (1.5pt) node[below=2pt] {$b$};
    \fill (2.6,1.62) circle (1.5pt) node[below=2pt] {$c$};
    \draw[<->,thick] (0.76,1.62) -- (1.44,1.62) node[midway,above=2pt] {$\sim$};
    \draw[->,thick] (1.76,1.62) -- (2.44,1.62) node[midway,above=2pt] {$\succ$};
    \draw[->,thick] (2.68,2.12) to[bend left=25] (0.52,2.12) node[midway,above=2pt] {$\succ$};
    \node[align=center,text width=3.0cm] at (1.6,0.52) {双向箭头读作无差异，\\ 单向粗箭头读作严格};
  \end{scope}
\end{tikzpicture}
\caption{二元关系的分解：无差异部分 $I(R)$ 与严格部分 $P(R)$}
\end{figure}
\figexplan{左面板把 $R\subseteq X\times X$ 画成单元格方阵：横轴为第一个元素、纵轴为第二个元素，落在 $R$ 内的格子涂深。若 $(x,y)$ 与 $(y,x)$ 同时落在 $R$ 内（深灰），该对元素无差异，属 $I(R)$；若只有 $(x,y)$ 落入（浅灰），则为严格偏好，属 $P(R)$。右面板把同一关系翻译成 $X=\{a,b,c\}$ 上的读法：双向箭头表示无差异、单向粗箭头表示严格偏好，并给出 $R=I(R)\cup P(R)$ 的分解恒等式。}

\begin{definition}[无差异部分与严格偏好部分]
对应于 $R\in\mathcal{B}$，其\textbf{对称部分}（无差异关系）与\textbf{非对称部分}（严格偏好关系）分别为
\begin{align}
I(R) &= \{(x,y)\in X\times X \mid (x,y)\in R \ \&\ (y,x)\in R\}, \label{eq:IR}\\
P(R) &= \{(x,y)\in X\times X \mid (x,y)\in R \ \&\ (y,x)\notin R\}. \label{eq:PR}
\end{align}
即：$x$ 与 $y$ 无差异 $\iff$ 双向弱偏好同时成立；$x$ 严格优于 $y$ $\iff$ 正向成立而反向不成立。
\end{definition}

\begin{remark}
这是本章一切定义的“语法底座”：后续所有公理、所有显示偏好关系，最终都归结为对
$R$、$I(R)$、$P(R)$ 的集合运算陈述。请熟练掌握 $(x,y)\in R$、$I(R)$、$P(R)$ 三者互推。
\end{remark}

\subsection{\texorpdfstring{$\succeq$/$\succ$/$\sim$}{≽/≻/∼} 记法与 $R$ 记法的互译}

上一小节用集合论语言定义了弱偏好关系 $R$ 及其分解；经济学文献与后续章节
习惯直接用关系符号书写。本章并行使用两套\textbf{完全等价}的记法：
§2.1--§2.3 以 $\succeq$ 记法为主，§2.4--§2.6 进入一般理论时以 $R$ 记法为主。

\begin{definition}[弱偏好、严格偏好、无差异]\label{def:prefsymbols}
设 $X$ 为所有可能选择的集合。
\begin{itemize}
  \item \textbf{弱偏好关系} $\succeq$：$X$ 上的二元关系，
  $x\succeq y$ 当且仅当“$x$ 至少和 $y$ 一样好”（at least as good as）；
  \item \textbf{严格偏好关系} $\succ$：$x\succ y$（“$x$ 严格优于 $y$”）当且仅当
  $x\succeq y$ 且 $y\not\succeq x$；
  \item \textbf{无差异关系} $\sim$：$x\sim y$（“二者无差异”）当且仅当
  $x\succeq y$ 且 $y\succeq x$。
\end{itemize}
\end{definition}

\textbf{互译词典}（务必熟练双向换算）：
\[
x\succeq y \iff (x,y)\in R;\qquad
x\succ y \iff (x,y)\in P(R);\qquad
x\sim y \iff (x,y)\in I(R).
\]

\textbf{符号总表}：
给定任意 $x,y\in X$，两个独立问题——“$x\succeq y$ 成立吗？”“$y\succeq x$ 成立吗？”——
各有是/否两答，组合出四种情形，恰好三种有名字：
\begin{center}
\small
\begin{tabular}{c|cc}
\toprule
 & $y\succeq x$ & $y\not\succeq x$ \\
\midrule
$x\succeq y$ & $x\sim y$（无差异） & $x\succ y$（严格偏好） \\
$x\not\succeq y$ & $y\succ x$ & 被\textbf{完备性}排除 \\
\bottomrule
\end{tabular}
\end{center}
因此，若 $\succeq$ 完备，可以把它看成一个\textbf{三值比较函数}
\[
\rho:\ X\times X\longrightarrow \{\succ,\ \sim,\ \prec\},
\]
即每对 $(x,y)$ 都被归类为三种关系之一——这正是“完备偏好 = 全盘可比”的精确含义。

\subsection{完备性、传递性与理性}

以下三个定义是全课的地基，是后续一切定理的输入。

\begin{definition}[完备性（completeness）]\label{def:complete}
$\succeq$ 在 $X$ 上\textbf{完备}，当且仅当
\[
\forall x,y\in X:\quad x\succeq y\ \text{或}\ y\succeq x.
\]
\end{definition}

\textbf{解读：}
\begin{enumerate}
  \item “或”是\textbf{相容析取}：$x\succeq y$ 与 $y\succeq x$ 允许同时成立——
  那正是无差异 $x\sim y$。完备性不排除“一样好”，只排除“比不了”。
  \item 被排除的情形\textbf{恰好一种}：$x\not\succeq y$ 且 $y\not\succeq x$
  ——即 2×2 总表第四象限的“不可比”。完备性的全部内容就是
  \textbf{禁止犹豫}：面对任何两个选项，必须能表态（哪怕表态为“无差异”）。
  \item \textbf{完备性 $\Rightarrow$ 自反性}：取 $y=x$，定义给出 $x\succeq x$ 或
  $x\succeq x$，恒真。
  反之不然：自反性只说每个元素与自身的比较，对不同元素之间不置一词。
  \item \textbf{经济内容：这是很强的行为假设。}现实中“合理的不可比”比比皆是：
  \begin{itemize}
    \item \textbf{Pareto 支配}：$x\ge y$ 且 $x\neq y$（每种商品都不少）是
    $\R^l$ 上的天然偏好，但两个“各有千秋”的消费束不可比——它不完备；
    \item \textbf{多维权衡}：学术职位 vs 业界职位，收入、自由、声望、风险
    权重难定，“说不清哪个好”是常态而非病态；
    \item \textbf{Bewley (1986)}：在 Knight 不确定性下，无法形成概率判断的
    决策者\textbf{理应}偏好不完备——“不表态”本身是对无知的诚实。
    Buridan 之驴（两堆完全相同的干草之间饿死）是文学化的反例。
  \end{itemize}
  \item 常见误解：完备性\textbf{不是}“对一切都有明确数值评价”，
  也不是“喜欢所有东西”——它只要求\textbf{两两之间必有一票}。
\end{enumerate}

\begin{definition}[传递性（transitivity）]\label{def:transitive}
$\succeq$ 在 $X$ 上\textbf{传递}，当且仅当
\[
\forall x,y,z\in X:\quad x\succeq y\ \text{且}\ y\succeq z\ \Rightarrow\ x\succeq z.
\]
\end{definition}

\textbf{解读：}
\begin{enumerate}
  \item 传递性\textbf{只}排除含严格环节的偏好环——
  无差异环 $x\sim y\sim z\sim x$ 无害（大家一样好，随便选），
  严格环 $x\succ y\succ z\succ x$ 致命（Alice 例 \ref{ex:alice}：
  每个选项都被另一个击败，$C(X,\succeq)=\varnothing$）。
  \item \textbf{货币泵论证（money pump）}——传递性的市场辩护：
  设某人有严格环 $x\succ y\succ z\succ x$，且偏好不随小额金钱改变。
  他手中持有 $x$。经纪人上门，利用 $z\succ x$ 提议：
  “付我 1 元，把你手里的 $x$ 换成你更喜欢的 $z$”——他愿意。
  再利用 $y\succ z$：“付 1 元，把 $z$ 换成你更喜欢的 $y$”——又愿意。
  再利用 $x\succ y$：“付 1 元，把 $y$ 换成你更喜欢的 $x$”——还是愿意。
  三轮过后他\textbf{回到起点}（手中仍是 $x$），净付 3 元，
  且同样的三轮可以无限重复——严格偏好环就是一棵
  可以被反复收割的摇钱树。“非理性偏好者在市场上活不久”，
  这是传递性的\textbf{演化辩护}：不满足它的行为主体会被系统性套干。
  \item \textbf{实验证据}：Tversky (1969, \emph{Psychological Review})
  报告了系统性的现场不传递（被试在成对比较中持续走出环）。
  传递性因而是\textbf{理想理性的假设}而非\textbf{真实行为的描述}——
  这正是 §2.1.4 弱化条件（拟传递、Suzumura 一致、无环）与
  行为经济学批判的入口。
  \item \textbf{计算的通行证}：传递性使“排序”良定义——$n$ 个选项只需
  $O(n\log n)$ 次两两比较即可排出全序（任何排序算法都依赖传递性）；
  不传递时“最好的那个”概念本身崩塌，选择规则可以为空
  （定理 \ref{thm:nonempty} 的归纳证明在情形 2 用了传递性，不可省）。
\end{enumerate}

\begin{definition}[理性偏好（rationality，MWG 定义 1.B.1）]\label{def:rational}
$\succeq$ 在 $X$ 上\textbf{理性}，当且仅当它\textbf{完备且传递}。
\end{definition}

\begin{remark}[“理性”是技术术语，不是价值判断]
\begin{itemize}
  \item 理性 $\neq$ 自私、$\neq$ 聪明、$\neq$ 目标高尚——
  它只要求偏好\textbf{结构自洽}，对\textbf{内容}不置一词。
  一个疯子完全一致的偏好也是“理性”的；一个善人犹豫不决的偏好不是。
  \item 反过来，两条公理也\textbf{不保证}行为合乎常识——
  它们管的是“账本自洽”，不管“账本上记了什么”。
  \item 但仅有这两条公理，已足够撑起本章全部理论：
  非空性定理 \ref{thm:nonempty}、效用表示（§2.2）、
  显示偏好刻画（§2.4--§2.6）的输入\textbf{只有}完备性与传递性。
  全部后续工作（§2.1.4 与 §2.6.4）都是在问：\textbf{这两条公理能弱化到什么程度，
  而理论大厦不塌？}
\end{itemize}
\end{remark}

\begin{figure}[H]
\centering
\begin{tikzpicture}[font=\small,
  pt/.style={circle, draw, inner sep=3pt},
  arr/.style={-{Stealth[length=2.2mm]}, thick}]
% 左：严格环
\begin{scope}
\node[pt] (x) at (0,1.5) {$x$};
\node[pt] (y) at (-1.5,0) {$y$};
\node[pt] (z) at (1.5,0) {$z$};
\draw[arr] (x) -- node[above left=1pt, font=\footnotesize]{$\succ$} (y);
\draw[arr] (y) -- node[below=3pt, font=\footnotesize]{$\succ$} (z);
\draw[arr] (z) -- node[above right=1pt, font=\footnotesize]{$\succ$} (x);
\node[font=\footnotesize\bfseries] at (0,-1.1) {严格环};
\end{scope}
% 右：传递序
\begin{scope}[shift={(6,0)}]
\node[pt] (a) at (0,1.5) {$x$};
\node[pt] (b) at (0,0) {$y$};
\node[pt] (c) at (0,-1.5) {$z$};
\draw[arr] (a) -- node[right=4pt, font=\footnotesize]{$\succ$} (b);
\draw[arr] (b) -- node[right=4pt, font=\footnotesize]{$\succ$} (c);
\draw[arr, dashed] (a) to[bend left=40] node[left=4pt, font=\footnotesize]{$\succ$} (c);
\node[font=\footnotesize\bfseries] at (0,-2.4) {传递序};
\end{scope}
\end{tikzpicture}
\caption{传递性在选择中的作用。左：严格环 $x\succ y\succ z\succ x$，
每个选项都被另一个严格击败，$C(X,\succeq)=\varnothing$，无人能被选中（例 \ref{ex:alice}）；
右：传递序 $x\succ y\succ z$，顶端元素是所有比较的胜利者，
$C(X,\succeq)=\{x\}$（定理 \ref{thm:nonempty}）。虚线是传递性“免费赠送”的比较：
$n$ 个元素只需 $n-1$ 次相邻比较即可定序。}
\label{fig:cyclevschain}
\end{figure}
\figexplan{左图为严格环 $x\succ y\succ z\succ x$：每个选项都被另一个严格击败，因此不存在最大元，$C(X,\succeq)=\varnothing$——无人能被选中。右图为传递序 $x\succ y\succ z$：顶端元素 $x$ 在所有比较中获胜，$C(X,\succeq)=\{x\}$。右侧的虚线弧是传递性“免费赠送”的比较：只需 $n-1$ 次相邻比较即可为 $n$ 个元素定序。}

\begin{proposition}[理性偏好的导出性质]\label{prop:derived}
设 $\succeq$ 理性（完备 + 传递）。则对所有 $x,y,z\in X$：
\begin{enumerate}
  \item $\succeq$ \textbf{自反}：$x\succeq x$；
  \item $\sim$ 是\textbf{等价关系}：自反（$x\sim x$）、对称（$x\sim y\Rightarrow y\sim x$）、
  传递（$x\sim y\sim z\Rightarrow x\sim z$）；
  \item $\succ$ \textbf{反自反}（$x\not\succ x$）且\textbf{传递}
  （$x\succ y\succ z\Rightarrow x\succ z$）；
  \item \textbf{混合蕴含}：$(x\succ y)\ \&\ (y\succeq z)\Rightarrow x\succ z$，且
  $(x\succeq y)\ \&\ (y\succ z)\Rightarrow x\succ z$。
\end{enumerate}
\end{proposition}

\begin{proof}
(1) 完备性用于 $(x,x)$：$x\succeq x$ 或 $x\succeq x$，恒成立。

(2) 自反性由 (1)；对称性即 $\sim$ 的定义；传递性：$x\sim y$ 给出
$x\succeq y$、$y\succeq x$；$y\sim z$ 给出 $y\succeq z$、$z\succeq y$；
由 $\succeq$ 传递，$x\succeq z$ 且 $z\succeq x$，即 $x\sim z$。

(3) 反自反：$x\succ x$ 要求 $x\not\succeq x$，与 (1) 矛盾。
传递：$x\succ y\succ z$ 给出 $x\succeq y\succeq z$，由 $\succeq$ 传递得
$x\succeq z$；若 $z\succeq x$，则 $z\succeq x\succeq y$ 给出 $z\succeq y$，
与 $y\succ z$ 矛盾。故 $x\succ z$。

(4) 只证第一式（第二式对称）。设 $x\succ y$ 且 $y\succeq z$。
由 $x\succeq y$、$y\succeq z$ 及传递性得 $x\succeq z$。
反设 $z\succeq x$：则 $y\succeq z$（已知）与 $z\succeq x$（反设）经传递得
$y\succeq x$，与 $x\succ y$（即 $y\not\succeq x$）矛盾。
故 $z\not\succeq x$，即 $x\succ z$。
\end{proof}

\begin{example}[Alice 的早餐：传递性为什么会破]\label{ex:alice}
$X=\{$热干面，豆皮，小笼包$\}$。Alice 的偏好为：
\[
\text{热干面}\succ\text{豆皮},\qquad
\text{豆皮}\succ\text{小笼包},\qquad
\text{小笼包}\succeq\text{热干面}.
\]
前两条由命题 \ref{prop:derived}(3) 推出 $\text{热干面}\succ\text{小笼包}$，
与第三条（小笼包 $\succeq$ 热干面）矛盾——Alice 违反的是\textbf{传递性}
（完备性三条比较均已给出，并未违反）。
后果是灾难性的：面对 $X$ 她\textbf{无法做出选择}——每个方案都被另一个严格击败，
$C(X,\succeq)=\varnothing$（见 §2.3.3 的非空性定理：传递性正是选择集非空的保证）。
课堂提问：若只保留前两条，Alice 的偏好完备吗？（不完备——缺热干面与小笼包的比较；
但此时 $C(X,\succeq)=\{$热干面$\}$ 仍然非空。）
\end{example}

\subsection{标准性质清单}

\begin{definition}[二元关系的标准性质]
对 $R\in\mathcal{B}$：
\begin{enumerate}
  \item[(a)] \textbf{完备性（completeness）}：$\forall x,y\in X,\ x\neq y \Rightarrow (x,y)\in R \vee (y,x)\in R$；
  \item[(b)] \textbf{自反性（reflexivity）}：$\forall x\in X:\ (x,x)\in R$；
  \item[(c)] \textbf{反自反性（irreflexivity）}：$\forall x\in X:\ (x,x)\notin R$；
  \item[(d)] \textbf{传递性（transitivity）}：$\forall x,y,z\in X:\ [(x,y)\in R \ \&\ (y,z)\in R] \Rightarrow (x,z)\in R$；
  \item[(e)] \textbf{对称性（symmetry）}：$(x,y)\in R \Rightarrow (y,x)\in R$；
  \item[(f)] \textbf{非对称性（asymmetry）}：$(x,y)\in R \Rightarrow (y,x)\notin R$；
  \item[(g)] \textbf{等价关系（equivalence）}：满足 (b)(d)(e)；
  \item[(h)] \textbf{序（ordering）}：满足 (a)(b)(d)；
  \item[(i)] \textbf{拟序（quasi ordering）}：满足 (b)(d)。
\end{enumerate}
\end{definition}

\noindent\textbf{课堂要点。}经济学里“理性偏好”的标准定义就是 (h)：完备 + 自反 + 传递。
完备性要求任意两者可比（不允许“我不知道选哪个”），传递性排除循环偏好。
传递性是整套理论的“承重墙”，但行为实验反复表明真实个体常违反它（Luce 的咖啡加糖例）——
这促使我们引入下一小节\textbf{弱于传递}的一致性概念。

\subsection{弱于传递的一致性：拟传递、Suzumura 一致、无环}

\begin{definition}[三种弱化的一致性条件]\label{def:weakcoherence}
对 $R\in\mathcal{B}$：
\begin{enumerate}
  \item[(j)] \textbf{拟传递性（quasi transitivity，Sen 1969）}：只要求\textbf{严格部分}传递：
  \[
    \forall x,y,z\in X:\ (x,y)\in P(R)\ \&\ (y,z)\in P(R)\ \Rightarrow\ (x,z)\in P(R).
  \]
  \item[(k)] \textbf{Suzumura 一致性（Suzumura consistency，Suzumura 1976a）}：
  \[
    \forall x,y\in X:\ (x,y)\in T(R)\ \Rightarrow\ (y,x)\notin P(R),
  \]
  其中 $T(R)$ 为 $R$ 的传递闭包（定义 \ref{def:TC}）。
  \item[(l)] \textbf{无环性（acyclicity）}：$\forall x\in X:\ (x,x)\notin T(P(R))$，
  即不存在严格偏好环。
\end{enumerate}
\end{definition}

\begin{definition}[传递闭包]\label{def:TC}
对 $R\in\mathcal{B}$，其\textbf{传递闭包} $T(R)$ 定义为：$(x,y)\in T(R)$ 当且仅当存在自然数
$t\ (2\le t<+\infty)$ 与序列 $x^0,x^1,\dots,x^t\in X$，满足 $x^0=x$，$(x^\tau,x^{\tau+1})\in R$
（$\tau=0,1,\dots,t-1$），且 $x^t=y$。
\end{definition}

可以验证算子 $T$ 是 $\mathcal{B}$ 上的\textbf{闭包运算}（Berge 1963）：
\[
\text{(T}_1\text{)}\ T(R)\supseteq R;\quad
\text{(T}_2\text{)}\ R^1\supseteq R^2 \Rightarrow T(R^1)\supseteq T(R^2);\quad
\text{(T}_3\text{)}\ T[T(R)]=T(R);\quad
\text{(T}_4\text{)}\ T(\varnothing)=\varnothing.
\]

\begin{theorem}[Suzumura 1983, p.\,12；Theorem I.1]\label{thm:TC}
对任意 $R\in\mathcal{B}$，$T(R)$ 是\textbf{包含 $R$ 的最小传递超集}。
\end{theorem}

\begin{proof}[完整证明（教材略去，此处补全）]
需证三件事：(i) $T(R)$ 是传递的；(ii) $R\subseteq T(R)$；(iii) 若 $Q\in\mathcal{B}$
传递且 $R\subseteq Q$，则 $T(R)\subseteq Q$。

\textbf{(i) 传递性。}设 $(x,y)\in T(R)$、$(y,z)\in T(R)$。由定义 \ref{def:TC}，
存在链 $x=x^0,x^1,\dots,x^t=y$ 与链 $y=y^0,y^1,\dots,y^s=z$，相邻项均属 $R$。
把两链首尾相接：
\[
x=x^0,\dots,x^t=y=y^0,y^1,\dots,y^s=z
\]
是一条从 $x$ 到 $z$ 的 $R$-链（长度 $t+s\ge 2$，因 $t,s\ge 2$……若允许单节链即
$t\ge 1$，结论相同），故 $(x,z)\in T(R)$。

\textbf{(ii) 包含关系。}$(x,y)\in R$ 时取单节链 $x^0=x,\ x^1=y$（$t=1$），
即得 $(x,y)\in T(R)$，故 $R\subseteq T(R)$。

\textbf{(iii) 最小性。}设 $Q$ 传递且 $R\subseteq Q$。任取 $(x,y)\in T(R)$，
存在 $R$-链 $x=x^0,\dots,x^t=y$。由 $R\subseteq Q$，链上每节
$(x^\tau,x^{\tau+1})\in Q$；对链长作归纳：$t=1$ 时 $(x,y)\in Q$ 显然；
设 $t\le k$ 时结论成立，$t=k+1$ 时由归纳假设 $(x,x^k)\in Q$、又 $(x^k,y)\in Q$，
由 $Q$ 的传递性得 $(x,y)\in Q$。故 $T(R)\subseteq Q$。

由 (iii) 取 $Q=T(R)$（由 (i) 它传递、由 (ii) 它含 $R$）知 $T(R)$ 自身是候选之一，
且它含于一切候选，故为\textbf{最小}传递超集。\textbf{推论}：$R$ 传递
$\iff T(R)=R$。
\end{proof}

\begin{remark}[Suzumura 一致性的直观]
条件 (k) 排除如下结构：存在有限链 $x^0=x,\ (x^\tau,x^{\tau+1})\in R,\ x^t=y$ 同时 $(y,x)\in P(R)$。
换言之，它同时排除\textbf{弱货币泵}与\textbf{强货币泵}：沿弱偏好链从 $x$ 走到 $y$ 之后，
不允许 $y$ 又严格优于 $x$。它比传递性弱（不要求链条闭合后仍有弱偏好），但足以阻止“被来回收割”。
\end{remark}

三种弱化概念与传递性的逻辑关系如下图（单向箭头不可反转；Suzumura 一致与拟传递\textbf{逻辑独立}）：

\begin{figure}[H]
\centering
\begin{tikzpicture}[
  node distance=1.1cm and 2.2cm,
  box/.style={draw, rounded corners=2pt, inner sep=6pt, font=\small},
  arr/.style={-{Stealth[length=2.5mm]}, thick}
]
\node[box] (tr) {传递性 Transitivity};
\node[box, right=of tr] (qt) {拟传递 Quasi transitivity};
\node[box, below=of tr] (sz) {Suzumura 一致性};
\node[box, right=of sz] (ac) {无环性 Acyclicity};
\draw[arr] (tr) -- (qt);
\draw[arr] (tr) -- (sz);
\draw[arr] (qt) -- (ac);
\draw[arr] (sz) -- (ac);
\end{tikzpicture}
\caption{一致性条件的强弱：传递性 $\Rightarrow$ 拟传递 $\Rightarrow$ 无环；传递性 $\Rightarrow$ Suzumura 一致 $\Rightarrow$ 无环。Suzumura 一致与拟传递互不蕴含（Suzumura 1983, p.\,244）。}
\label{fig:coherence}
\end{figure}
\figexplan{箭头表示不可反转的逻辑蕴含：传递性同时推出拟传递与 Suzumura 一致性，两者再各自推出无环性。拟传递与 Suzumura 一致性位于图的中间两角且互不相连，表示二者逻辑独立、互不蕴含——这正是弱化处理最需要留意的地方。}

\begin{example}[Luce 1956：咖啡加糖——无差异为何不该传递]\label{ex:luce}
设 $X=\{x_0,x_1,\dots,x_n\}$，$x_k$ 是一杯加了 $k$ 粒糖的咖啡。人的味觉存在
\textbf{最小可觉差}（just noticeable difference, JND）：相邻两杯分辨不出，
$x_k\sim x_{k+1}$（$k=0,\dots,n-1$）；但 $0$ 粒糖与 $n$ 粒糖（$n$ 很大）
一苦一甜，$x_n\succ x_0$。
\begin{itemize}
  \item 若要求 $I(R)$ 传递：$x_0\sim x_1\sim\cdots\sim x_n$ 推出 $x_0\sim x_n$，矛盾；
  \item 若只要求\textbf{拟传递}：严格部分 $x_n\succ x_0$ 与任何已有严格偏好不冲突，
  无任何矛盾。
\end{itemize}
这正是 Sen (1969) 主张以拟传递替代完全传递的动机之一：\textbf{无差异的不可分辨性
（imperceptible differences）使 $I(R)$ 天然可能不传递}，而严格偏好 $P(R)$
基于可分辨的差异，传递性仍可辩护。Luce (1956) 由此发展出\textbf{半序}
（semiorder）理论：$(x,y)\in P(R)$ 当且仅当 $u(x)>u(y)+\delta$（$\delta>0$ 为
感觉阈限），其无差异类恰是“叠合区间”——拟传递而不传递。
同一思想可归一化为“一单位 JND”：$x\succeq y \iff u(x)\ge u(y)-1$
——此时 $\succ$ 传递（$x\succ y \iff u(x)>u(y)+1$，即优势超过阈限，
“超过阈限的差距”当然可传递）而 $\succeq$ 不传递：取 $u(x)=2$、$u(y)=2.9$、$u(z)=3.8$，
则 $x\succeq y$（$2\ge 1.9$）、$y\succeq z$（$2.9\ge 2.8$），但
$x\not\succeq z$（$2<2.8$）——弱偏好链在阈值处断裂。
\end{example}

\begin{figure}[H]
\centering
\begin{tikzpicture}[font=\scriptsize,scale=1.12,transform shape]
  \draw[thick,fill=gray!15] (0,0) rectangle (3.1,2.9);
  \node[font=\small\bfseries] at (1.55,2.56) {(a) 含环的原偏好};
  \fill (1.55,2.00) circle (1.5pt) node[above=2pt] {$a$};
  \fill (0.65,0.95) circle (1.5pt) node[below=2pt] {$b$};
  \fill (2.45,0.95) circle (1.5pt) node[below=2pt] {$c$};
  \draw[->,thick] (1.30,1.82) -- (0.85,1.15);
  \draw[->,thick] (1.05,0.90) -- (2.25,0.90);
  \draw[->,thick] (2.50,1.15) -- (1.85,1.80);
  \node[align=center,fill=white,inner sep=2pt,text width=2.65cm] at (1.55,0.32) {$a\succ b\succ c\succ a$\\[-1pt]非传递};
  \begin{scope}[xshift=6.2cm]
    \draw[thick,fill=gray!15] (0,0) rectangle (4.4,2.9);
    \node[font=\small\bfseries] at (2.2,2.56) {(b) 传递扩张：无穷序列};
    \foreach \i in {0,1,2,3} \fill (0.5+0.9*\i,1.75) circle (1.5pt);
    \node[below=2pt] at (0.5,1.75) {$a_1$};
    \node[below=2pt] at (1.4,1.75) {$b_1$};
    \node[below=2pt] at (2.3,1.75) {$c_1$};
    \node[below=2pt] at (3.2,1.75) {$a_2$};
    \node[above=2pt] at (0.5,1.75) {$a$};
    \node[above=2pt] at (1.4,1.75) {$b$};
    \node[above=2pt] at (2.3,1.75) {$c$};
    \node[above=2pt] at (3.2,1.75) {$a$};
    \node at (3.75,1.75) {$\cdots$};
    \draw[->,thick] (0.62,1.45) -- (1.28,1.45);
    \draw[->,thick] (1.52,1.45) -- (2.18,1.45);
    \draw[->,thick] (2.42,1.45) -- (3.08,1.45);
    \node[align=center,fill=white,inner sep=2pt,text width=3.8cm] at (2.2,0.72) {无穷多个“复制点”按序排定，\\原偏好嵌入其中且被保留};
    \node[align=center,fill=white,inner sep=2pt,text width=3.8cm] at (2.2,0.24) {扩张后的关系是传递的\\（原关系是其子关系）};
  \end{scope}
\end{tikzpicture}
\caption{Szpilrajn 扩张定理：把含环偏好嵌入传递关系}
\end{figure}
\figexplan{左面板是一个含三环的偏好 $a\succ b\succ c\succ a$，它满足完备性但不满足传递性，因而无法被任何传递关系（特别是效用）表示。右面板给出 Szpilrajn 扩张的做法：把每个元素复制成无穷多个副本 $a_1,b_1,c_1,a_2,\dots$，按线性顺序排定，使扩张后的关系可传递；同时通过“同序映射”保证原偏好被原样保留，于是原关系成为传递关系的子关系。}

\subsection{扩张与 Szpilrajn 定理}

\begin{definition}[扩张与子关系]
设 $R, R^*\in\mathcal{B}$。若 $R\subseteq R^*$ 且 $P(R)\subseteq P(R^*)$，称 $R^*$ 是 $R$
的一个\textbf{扩张（extension）}，$R$ 是 $R^*$ 的\textbf{子关系（subrelation）}。
若 $R^*$ 还是一个序，称其为 $R$ 的\textbf{序扩张（ordering extension）}。
\end{definition}

\begin{theorem}[Szpilrajn 1930]\label{thm:szpilrajn}
任何拟序 $R$（自反 + 传递）都存在一个序扩张 $R^*$。
\end{theorem}

\begin{remark}
Szpilrajn 定理是本章的“隐藏支柱”：它保证\textbf{部分可比}的偏好总可以\textbf{补全}为完备序而不反转任何已有严格比较。
在显示偏好理论中，我们先用选择数据构造出一个（可能不完备的）显示偏好关系，再靠扩张定理把它“补齐”为理性化所需的序。
\end{remark}

\begin{proof}[定理 \ref{thm:szpilrajn} 的证明思路（Zorn 引理）]
证明分两步：先证\textbf{单点补全引理}，再用极大原理收尾。

\textbf{第一步（关键引理）。}设 $R$ 是拟序，$x,y\in X$ 且 $(x,y)\notin R$、
$(y,x)\notin R$（$x$ 与 $y$ 不可比）。定义
\[
R^1 = R\ \cup\ \{(x,y)\}\ \cup\ \{(z,w)\in X\times X \mid (z,x)\in R\ \&\ (y,w)\in R\},
\]
即：强行宣布 $x\succeq y$，并把“所有不弱于 $x$ 的 $z$”与“所有不优于 $y$ 的 $w$”
之间也补上 $z\succeq w$。可逐条验证 $R^1$ 仍自反、传递：
若 $(a,b),(b,c)\in R^1$，按各项来源（旧 $R$ 或新补的）分四种情形，
用 $R$ 的传递性逐一合并，均得 $(a,c)\in R^1$；自反性显然。
并且 $P(R)\subseteq P(R^1)$：原有的严格比较不会被反转
（若 $(a,b)\in P(R)$ 而 $(b,a)\in R^1$ 为新增，则 $(b,x),(y,a)\in R$，
结合 $(x,y)$ 新增前的结构与 $R$ 传递会推出 $(y,x)\in R$，与不可比假设矛盾）。
故 $R^1$ 是 $R$ 的真扩张，且 $x,y$ 在 $R^1$ 下可比。

\textbf{第二步（Zorn 引理）。}令
$\mathfrak{Q}=\{Q\in\mathcal{B}\mid Q \text{ 拟序},\ R\subseteq Q,\ P(R)\subseteq P(Q)\}$，
以集合包含为序。任取链 $\{Q_\lambda\}\subseteq\mathfrak{Q}$，其并
$\bigcup_\lambda Q_\lambda$ 仍自反（各 $Q_\lambda$ 均含对角线）、仍传递
（$(a,b)\in Q_{\lambda_1}$、$(b,c)\in Q_{\lambda_2}$ 时取
$\lambda=\max\{\lambda_1,\lambda_2\}$，两项同处一个 $Q_\lambda$ 中），
且保持 $P(R)$，故并属于 $\mathfrak{Q}$，是该链的上界。
由 Zorn 引理，$\mathfrak{Q}$ 存在极大元 $R^*$。
若 $R^*$ 不完备，则有不可比对 $x,y$，第一步给出真扩张 $R^1\supsetneq R^*$，
与极大性矛盾。故 $R^*$ 完备，即为所求序扩张。
\end{proof}

\begin{remark}[构造性与非构造性]
$X$ 有限时无需 Zorn 引理：不可比对至多有限，反复用第一步引理（至多
$\binom{|X|}{2}$ 次）即可补全——这给出 Szpilrajn 定理的\textbf{有限构造性证明}
（课堂可演算）。无限情形必须借助选择公理的等价形式，说明“完备化”本质上是一个
\textbf{非构造}操作：序扩张存在，但没有算法能把它写出来。
\end{remark}

% ============================================================
\section{效用表示}
% ============================================================

\subsection{表示的定义与有限集表示定理}

\begin{definition}[效用表示]
效用函数 $u: X\to\R$ \textbf{表示}偏好关系 $\succeq$，若对所有 $x,y\in X$：
\[
x\succeq y \iff u(x)\ge u(y).
\]
\end{definition}

\begin{theorem}[有限集效用表示定理]\label{thm:finiteutility}
有限集 $X$ 上任何\textbf{完备且传递}的偏好关系 $\succeq$，都可被某个效用函数
$u: X\to\{1,\dots,n\}$ 表示，其中 $n\equiv |X|$。
\end{theorem}

\noindent\textbf{直觉论证}（“逐层剥皮法”）：
\begin{enumerate}
  \item 把 $X$ 的“顶层”元素（选择集 $C(X,\succeq)$ 中的元素）赋予效用 $n=|X|$；
  \item 丢弃它们，剩下集合 $X'$；
  \item 若 $X'=\varnothing$ 则结束，否则对 $X'$ 重复第 1 步。
\end{enumerate}
完备 + 传递保证每一步“顶层”非空，有限性保证过程终止。

\begin{proof}[定理 \ref{thm:finiteutility} 的严格证明（对 $n\equiv|X|$ 作数学归纳）]
$n=0$ 时 $X=\varnothing$，定理平凡成立。

假设定理对至多 $n$ 个元素的集合成立。考虑 $|X|=n+1$。
完备传递偏好下选择集 $C(X,\succeq)\neq\varnothing$，故
$Y\equiv X\setminus C(X,\succeq)$ 至多有 $n$ 个元素。
由归纳假设，$Y$ 上的偏好可被某个 $u: Y\to\{1,\dots,n\}$ 表示。
把 $u$ 扩张到 $X$：对所有 $x\in C(X,\succeq)$ 令 $u(x)=n+1$。

下面验证扩张后的 $u$ 在 $X$ 上表示 $\succeq$，即对所有 $x,y\in X$：
$x\succeq y \iff u(x)\ge u(y)$。

\textbf{（$\Rightarrow$）}设 $x\succeq y$。
若 $x\in C(X,\succeq)$，则 $u(x)=n+1\ge u(y)$。
若 $x\notin C(X,\succeq)$，则 $y\notin C(X,\succeq)$（否则 $y$ 弱优于一切包括 $x$，
与 $x\succeq y$ 及 $x$ 不在选择集矛盾——严格说：$y\in C(X,\succeq)$ 意味着
$y\succeq x$，结合 $x\succeq y$ 得 $x\sim y$，而 $x$ 应与 $y$ 同在选择集中）。
于是 $x,y\in Y$，由 $u$ 在 $Y$ 上表示 $\succeq$（及传递性），$u(x)\ge u(y)$。

\textbf{（$\Leftarrow$）}设 $u(x)\ge u(y)$。
若 $u(x)=n+1$，则 $x\in C(X,\succeq)$，故 $x\succeq y$。
若 $u(x)\le n$，则 $x,y\in Y$，由 $u$ 在 $Y$ 上表示 $\succeq$ 得 $x\succeq y$。

归纳成立，定理得证。
\end{proof}

\begin{theorem}[效用表示蕴含理性]\label{thm:utilityrational}
若效用函数 $u: X\to\R$ 表示 $X$ 上的 $\succeq$，则 $\succeq$ 是理性的（完备 + 传递）。
\end{theorem}

\begin{proof}
\textbf{完备性}：对任意 $x,y\in X$，$u(x),u(y)\in\R$，实数全序给出 $u(x)\ge u(y)$ 或
$u(y)\ge u(x)$；由表示定义，$x\succeq y$ 或 $y\succeq x$。

\textbf{传递性}：设 $x\succeq y\succeq z$。由表示定义 $u(x)\ge u(y)\ge u(z)$，
实数传递性给出 $u(x)\ge u(z)$，故 $x\succeq z$。
\end{proof}

\begin{remark}
定理 \ref{thm:finiteutility} 与 \ref{thm:utilityrational} 合起来给出\textbf{有限集上的完全刻画}：
\[
\succeq \text{ 完备且传递} \iff \succeq \text{ 可被效用函数表示} \qquad (X \text{ 有限}).
\]
方向“$\Leftarrow$”永远成立（与基数无关）；方向“$\Rightarrow$”在无限集上\textbf{会失效}，见下。
\end{remark}

\begin{example}[可数集上的表示：枚举插入法]\label{ex:countable}
设 $X=\{x_1,x_2,\dots\}$ 可数，$\succeq$ 完备传递。归纳地构造 $u$：
已给前 $n$ 项赋好值且保持序关系，处理 $x_{n+1}$ 时分三种情形：
\begin{enumerate}
  \item 若 $x_{n+1}\succ x_i$ 对所有 $i\le n$：令 $u(x_{n+1})=\max_{i\le n}u(x_i)+1$；
  \item 若 $x_i\succ x_{n+1}$ 对所有 $i\le n$：令 $u(x_{n+1})=\min_{i\le n}u(x_i)-1$；
  \item 否则存在“夹住”它的相邻已赋值项 $a\succ x_{n+1}\succ b$
  （或与某项 $x_j$ 无差异），令 $u(x_{n+1})=\tfrac{1}{2}(u(a)+u(b))$
  （无差异情形令 $u(x_{n+1})=u(x_j)$）。
\end{enumerate}
每步插入保持已有序关系不变，故极限 $u:X\to\R$ 表示 $\succeq$。
\textbf{要点}：可数集的表示\textbf{不需要连续性假设}——归纳永远“有下一步”。
真正的障碍出现在\textbf{不可数}集：有理数网不足以标记不可数多个无差异类时
（词典序 $X=[0,1]^2$），表示才彻底失败（见下节反例）。
这也解释了 Cantor--Debreu 传统中“可数序稠密子集”条件为何是表示定理的核心。
\end{example}

\subsection{\texorpdfstring{当 $|X|=\infty$}{当 |X| 为无限集}：词典序偏好反例}

$X$ 无限时，归纳证明走不通，但表示\textbf{有时仍可能}：
\begin{example}[可表示的无限偏好]
$\R_+$ 上的偏好 $x_1\succeq x_2 \iff x_1\ge x_2$，由 $u(x)=x$ 表示（也可由任意单调变换表示）。
\end{example}

但无限时\textbf{并非总能}表示：

\begin{example}[词典序偏好（lexicographic preferences）不可表示]\label{ex:lexico}
$[0,1]^2\subseteq\R^2$ 上定义：$(x_1,y_1)\succeq(x_2,y_2)$ 当且仅当
\begin{itemize}
  \item $x_1>x_2$，或
  \item $x_1=x_2$ 且 $y_1\ge y_2$。
\end{itemize}
该偏好完备且传递，但\textbf{不能被任何效用函数表示}。
\end{example}

\begin{proof}[不可表示性的论证要点]
两条等价思路：
\begin{enumerate}
  \item \textbf{无差异曲线视角}：词典序偏好\textbf{没有非平凡的无差异曲线}——
  每个点的无差异集是单点。任何效用函数若要表示它，就必须把正方形中的\textbf{不可数多个}点
  一一映射到实数线上互不相同的值，且保持序。
  \item \textbf{序保持单射的不可能性}：这样的 $u$ 将是从 $[0,1]^2$ 到 $\R$ 的
  \textbf{保序一一映射}。但对每个 $x\in[0,1]$，点 $(x,0)$ 与 $(x,1)$ 满足
  $u(x,0)<u(x,1)$，从而每个 $x$ 对应一个\textbf{非退化区间} $(u(x,0),u(x,1))$，
  且不同 $x$ 的区间互不相交。实数线上互不相交的非退化开区间\textbf{至多可数个}
  （每个区间含一个有理数），与 $[0,1]$ 不可数矛盾。
\end{enumerate}
直观记忆：两种无穷“大小不同”——正方形里的序结构太“密”，塞不进一条线。
\end{proof}

\begin{remark}
无限集上何时\textbf{能}表示？需要额外的\textbf{可数序稠密}条件（Debreu 表示定理：
完备传递 + 连续 $\Rightarrow$ 存在连续效用表示）。本课程后面用到时会回到这一点；
这里只需记住：\textbf{“理性偏好 $\Rightarrow$ 效用表示”是有条件的，有限集免费，无限集要钱（连续性）。}
\end{remark}

\begin{figure}[H]
\centering
\begin{tikzpicture}[font=\scriptsize,scale=1.08,transform shape]
  \draw[thick,fill=gray!15] (0,0) rectangle (3.2,2.6);
  \node[font=\small\bfseries] at (1.6,2.28) {(a) 表示一：赋值 $(0,1,2,3)$};
  \draw[thick] (0.5,1.35) -- (2.7,1.35);
  \foreach \i in {0,1,2,3}{
    \fill (0.5+0.7*\i,1.35) circle (1.5pt);
    \node[below=3pt] at (0.5+0.7*\i,1.35) {$a_{\i}$};
    \node[above=3pt] at (0.5+0.7*\i,1.35) {$\i$};
  }
  \node[align=center] at (1.6,0.55) {任意严格递增变换 $f$ 都给出同一偏好};
  \draw[->,thick] (3.35,1.35) -- (4.15,1.35);
  \node[above=3pt] at (3.75,1.35) {$f$ 严格递增};
  \begin{scope}[xshift=7.0cm]
    \draw[thick,fill=gray!15] (0,0) rectangle (3.2,2.6);
    \node[font=\small\bfseries] at (1.6,2.28) {(b) 表示二：赋值 $(0,10,20,30)$};
    \draw[thick] (0.5,1.35) -- (2.7,1.35);
    \foreach \i/\v in {0/0,1/10,2/20,3/30}{
      \fill (0.5+0.7*\i,1.35) circle (1.5pt);
      \node[below=3pt] at (0.5+0.7*\i,1.35) {$a_{\i}$};
      \node[above=3pt] at (0.5+0.7*\i,1.35) {$\v$};
    }
    \node[align=center] at (1.6,0.55) {与 (a) 序相同，但数值本身无经济含义};
  \end{scope}
  \node[align=center,fill=white,inner sep=2pt,text width=9.0cm,font=\scriptsize] at (5.1,-0.72) {序数性：效用只编码偏好顺序，不编码“量”；人际比较无法从效用数值读出。};
\end{tikzpicture}
\caption{效用的序数性：表示不唯一，但序唯一}
\end{figure}
\figexplan{两个面板给出同一偏好 $a_0\succ a_1\succ a_2\succ a_3$ 的两个效用表示：$(0,1,2,3)$ 与 $(0,10,20,30)$。二者只相差一个严格递增变换 $f$，因此诱导完全相同的无差异类与选择。图形要点：效用的“水平”与“间距”可以任意改变，故只有序有意义；跨人的效用数值比较（人际比较）无法由偏好表示本身推出。}

\subsection{效用的序数性与人际比较的失败}

注意表示定义只要求
\[
x\succeq y \iff u(x)\ge u(y),
\]
因此 $u$ 的任意\textbf{严格递增变换} $f\circ u$ 同样表示 $\succeq$：
\begin{itemize}
  \item “表示 $\succeq$”这一性质是\textbf{序数的（ordinal）}；
  \item 根本不存在有意义的“效用单位”（there is no such thing as a ``util''）。
\end{itemize}

由此引出\textbf{人际比较的失败}：
\begin{enumerate}
  \item 仅凭选择理论无法比较张三的 $u(x)$ 与李四的 $u(x)$——两人的 $u$ 各自可乘任意正的仿射（甚至任意单调）变换，
  跨人比较没有锚点。这对\textbf{功利主义}（把各人效用加总成社会福利）的原始纲领是沉重打击。
  \item \textbf{Rawls}（上承 Kant 一线）的回应：放弃跨人比较，只要求我们站在\textbf{单一选择者}的位置
  （“无知之幕”后）做选择。
  \item 行为层面的修补：“\textbf{最小可觉差}（just noticeable difference）”建议把 1 个 util
  定义为个体能察觉的最小差异，即 $x\succeq y \iff u(x)\ge u(y)-1$。
  注意此时 $\succ$ 仍传递，但 $\succeq$ \textbf{不再传递}（$a$ 与 $b$ 无差异、$b$ 与 $c$ 无差异，
  但 $a$ 明显优于 $c$）——这从行为上解释了传递性为何可疑。
\end{enumerate}

% ============================================================
\section{显示偏好路径：从选择反推偏好}
% ============================================================

\subsection{方法论：实证起源、规范起源与四个追问}

\textbf{选择理论要回答什么？}这门理论瞄准两类问题：
\begin{description}
  \item[实证问题（positive）] 理解\textbf{个体自利如何驱动更大的经济系统}——
  价格形成、市场出清、比较静态，全部从“个体在约束下选最好”这一微观假设推出；
  \item[规范问题（normative）] 为\textbf{功利主义计算提供客观标准}——
  福利比较、政策评估需要一把“效用”尺子。
\end{description}

\textbf{两条建模路线的对照}：
\begin{center}
\small
\begin{tabular}{p{3.4cm}p{5.2cm}p{5.2cm}}
\toprule
 & \textbf{常规路线（preference-based）} & \textbf{显示偏好路线（choice-based）} \\
\midrule
起点 & 偏好 $\succeq$（在心中，不可观察） & 选择 $C_R(\cdot)$（在市场上，可观察） \\
终点 & 与偏好相容的选择 $C(B,\succeq)$ & 与选择相容的偏好 \\
核心问题 & “理性的人会怎么选？” & “这些选择像不像理性人选的？” \\
代表定理 & 非空性定理 \ref{thm:nonempty}、
效用表示（§2.2） & 枢纽定理 \ref{thm:harp}、
Afriat 定理（§2.7.5） \\
\bottomrule
\end{tabular}
\end{center}

\textbf{认识论的不对称}是理解全章的钥匙：
偏好和效用都在人脑里，原则上\textbf{不可观察}；
选择在市场上、问卷里、实验中，原则上\textbf{可观察}。
Samuelson 1938 年的纲领（操作主义在经济学中的落实）：
\emph{科学概念必须由可观察的操作来定义}——效用测不出来，选择测得出来，
因此理性选择理论的\textbf{全部可检验内容必须能翻译成选择语言}。
显示偏好方法就是这次翻译：从观察到的选择出发，
问什么样的偏好与之\textbf{相容（compatible）}。

由此引出四个层层递进的追问：

\begin{enumerate}
  \item \textbf{我们能否检验理性选择理论？怎么检验？}\\
  科学哲学的要求（Popper 可证伪性）：一个理论必须给出
  “\emph{若观察到这样的事件，理论就是错的}”的明确条件，否则不是科学。
  显示偏好理论的伟大之处正在于：检验条件（WARP/SARP/HARP/GARP）
  是\textbf{纯数据性质}——拿到选择数据即可逐条核验，
  不需要对不可观察的偏好作任何猜测。这是本章 §2.4--§2.6 的全部内容。
  \item \textbf{观察到的选择与“最大化某个目标函数”一致吗？}\\
  即\textbf{理性化的存在性}问题：是否存在某个完备传递的 $\succeq$ 使
  $C_R(\cdot)=C(\cdot,\succeq)$？§2.3.4 给出必要条件（显示偏好推断），
  §2.5 的枢纽定理给出充要条件（HARP）。
  \item \textbf{若一致，能否恢复出目标函数？}\\
  即\textbf{识别（identification）}问题。答案是“能，但不唯一”：
  理性化偏好一般不唯一——$R_C$ 只钉死被观察到的比较，
  未观察之处自由（§2.6 的反例集中，同一选择函数常有多个理性化）。
  恢复的标准步骤：$R_C$ $\to$ 取闭包 $\to$ Szpilrajn 扩张为序
  $\to$ Debreu 表示为 $u$（即 §2.7.1 Houthakker 三步证明的骨架）。
  \item \textbf{恢复出来后能做什么？人际比较行得通吗？}\\
  答案的一半是冷水：效用的\textbf{序数性}使“甲的 10 单位效用
  是否多于乙的 8 单位”没有意义（§2.2.3）——功利主义加总在
  本框架内\textbf{没有合法基础}，这是对规范起源的沉重打击。
  两种著名回应：Rawls（沿 Kant 传统）要求只考虑\textbf{单一选择者}
  在无知之幕后的判断，绕开人际比较；“最小可觉差”（JND）
  尝试以可分辨阈值定义效用单位（例 \ref{ex:luce}）。
  另一半是好消息：实证分析\textbf{不需要}人际比较——
  比较静态、Afriat 检验、福利的 Pareto 标准都只用量价数据。
\end{enumerate}

\subsection{选择函数与定义域}

\begin{definition}[选择函数与选择集]
设 $\mathcal{K}$ 是 $X$ 的一族非空子集（称为\textbf{定义域}）。$\mathcal{K}$ 上的
\textbf{选择函数} $C$ 把每个 $S\in\mathcal{K}$ 映到 $S$ 的一个非空子集 $C(S)\subseteq S$。
$S$ 称为\textbf{机会集（opportunity set）}，$C(S)$ 称为从 $S$ 中选出的\textbf{选择集}。
\end{definition}

三类常用定义域（理论适用范围的关键）：
\begin{itemize}
  \item $\mathcal{K}_F$：\textbf{有限域}——$X$ 的所有非空有限子集（个体选择与社会选择理论的标准设定）；
  \item $\mathcal{K}_{CB}$：\textbf{竞争预算域}——$X\subseteq\R_+^l$ 时所有预算集
  $B(\boldsymbol{p},M)=\{\boldsymbol{x}\in X\mid \boldsymbol{p}\boldsymbol{x}\le M\}$（传统消费者理论，Samuelson--Houthakker 的舞台）；
  \item $\mathcal{K}_A$：\textbf{任意域}——不加任何结构限制的非空集族（Suzumura 的一般理论）。
\end{itemize}

\begin{remark}
Arrow 的洞见：当只考虑预算集时，弱显示偏好公理\textbf{不足以}推出序的存在；
但当选择定义在\textbf{所有有限集}上时，弱公理就足够了。
“把选择函数的定义域扩展到所有有限集，需求函数视角将被大大简化”——
这直接催生了 Sen (1971)、Plott (1973)、Suzumura (1976) 等后续工作。
\end{remark}

\subsection{\texorpdfstring{选择规则 $C(B,\succeq)$}{选择规则 C(B,≽)}：从偏好到行为}

建模个体决策有\textbf{两条路线}：
\begin{enumerate}
  \item \textbf{常规路线（preference-based）}：从偏好出发，问什么选择与之相容；
  \item \textbf{显示偏好路线（choice-based）}：从观察到的选择出发，问什么偏好与之相容
  ——§2.3.4 起的主线。
\end{enumerate}
本小节先把常规路线走通。

\begin{definition}[选择规则]\label{def:choicerule}
给定 $X$ 上的偏好 $\succeq$ 与机会集（选择集）$B\subseteq X$，
\textbf{选择规则}是一个对应（correspondence），给出 $B$ 中所有“最好”的元素：
\begin{equation}
C(B,\succeq)\equiv\{x\in B:\ x\succeq y\ \text{对所有}\ y\in B\}. \label{eq:CB}
\end{equation}
\end{definition}

用 $R$ 记法翻译：$C(B,\succeq)$ 正是下节的 $R$-最优点集 $G(B,R)$
（定义见下）。选择理论的第一问是：\textbf{这个集合什么时候非空？}
（Alice 例 \ref{ex:alice} 已提示：传递性缺失时它可以为空。）

\begin{theorem}[选择规则的非空性（有限情形）]\label{thm:nonempty}
设 $\succeq$ 完备且传递，$B$ 有限非空。则 $C(B,\succeq)\neq\varnothing$。
\end{theorem}

\begin{proof}[对 $|B|$ 作数学归纳]
$|B|=1$：$B=\{x\}$，由完备性 $x\succeq x$（命题 \ref{prop:derived}(1)），
故 $x\in C(B,\succeq)$。

设对一切 $|B|=n\ge 1$ 结论成立。考虑 $|A|=n+1$，写 $A=B\cup\{x\}$（$|B|=n$）。
由归纳假设取 $y\in C(B,\succeq)$。由完备性只有两种情形：
\begin{enumerate}
  \item $y\succeq x$：$y$ 已不弱于 $B$ 中一切（归纳假设），又不弱于 $x$，
  故 $y\in C(A,\succeq)$；
  \item $x\succeq y$：$y$ 不弱于 $B$ 中一切，由\textbf{传递性} $x$ 不弱于
  $B$ 中一切，又 $x\succeq x$，故 $x\in C(A,\succeq)$。
\end{enumerate}
两情形均得 $C(A,\succeq)\neq\varnothing$，归纳成立。
\end{proof}

\begin{remark}[这个定理用足了什么]
基例用\textbf{自反}（来自完备性），归纳步用\textbf{完备性 + 传递性}——
理性偏好的两条公理各司其职，缺一不可：
缺完备性，$y$ 与 $x$ 可能不可比，两情形都不覆盖；
缺传递性，情形 2 的 $x$ 打不过 $B$ 中某个 $z$（Alice 例即是如此）。
$B$ 无限时结论失效，需要追加\textbf{拓扑条件}（如 $\succeq$ 连续 + $B$ 紧，
由 Weierstrass 定理保证最大值存在）——这是 §2.7.5 Debreu 连续性条件
的经济学用途之一。
\end{remark}

\subsection{最大点集与最优点集；理性化}

\begin{definition}[$R$-最大点集与 $R$-最优点集]
设 $R$ 为偏好关系，$S\subseteq X$：
\begin{align}
M(S,R) &= \{x\in X \mid x\in S \ \&\ (y,x)\notin P(R)\ \text{对所有}\ y\in S\},\\
G(S,R) &= \{x\in X \mid x\in S \ \&\ (x,y)\in R\ \text{对所有}\ y\in S\}.
\end{align}
$M(S,R)$ 是“没有谁能严格击败我”的集合（maximal），$G(S,R)$ 是“我不弱于任何人”的集合（greatest）。
\end{definition}

\begin{remark}\label{rem:GM}
(i) 恒有 $G(S,R)\subseteq M(S,R)$；(ii) 若 $R$ 完备，则 $G(S,R)=M(S,R)$。
对 $x,y\in G(S,R)$ 有 $(x,y)\in I_R$；对 $x,y\in M(S,R)$ 则要么 $(x,y)\in I_R$，
要么二者不可比（$(x,y)\notin R$ 且 $(y,x)\notin R$）。
\end{remark}

\begin{definition}[理性化、G-理性、M-理性]\label{def:rationalize}
$\mathcal{K}$ 上的选择函数 $C$ 称为\textbf{理性的（rational）}，若存在二元关系 $R\in\mathcal{B}$ 使
\[
\forall S\in\mathcal{K}:\ C(S)=G(S,R)=\{x^*\in S\mid \forall x\in S:\ (x^*,x)\in R\}.
\]
这样的 $R$ 称为 $C$ 的一个\textbf{理性化（rationalization）}。
更精细地：
\begin{itemize}
  \item[(a)] $C$ 是 \textbf{G-理性}的，若存在 $R$ 使 $C(S)=G(S,R)$ 对所有 $S$ 成立；
  \item[(b)] $C$ 是 \textbf{M-理性}的，若存在 $R$ 使 $C(S)=M(S,R)$ 对所有 $S$ 成立。
\end{itemize}
若 $R$ 还是完备的（完备的 + 传递的），相应称 $C$ 为完备理性（完全理性 / full rational）的。
\end{definition}

\begin{remark}[M-理性 $\Rightarrow$ G-理性，反之不然]\label{rem:MG}
设 $C$ 以 $R$ 为理性化是 M-理性的。定义 $R'$：$(x,y)\in R' \iff (y,x)\notin P_R$。
则 $R'$ 完备，且 $M(S,R)=G(S,R')$ 对所有 $S$ 成立，故 $C$ 是 G-理性的。
反向不成立：存在 $X=\{x,y,z\}$ 上的例子（$S_1=\{x,y\},S_2=\{x,z\},S_3=X$，
$C(S_1)=S_1,C(S_2)=S_2,C(S_3)=\{x\}$）是 G-理性而非 M-理性的（Suzumura 书 p.\,118--119 的构造）。
直观地说，G-理性允许“不可比”作为不选的理由，M-理性不允许，因此 M-理性更强。
\end{remark}

\subsection{理性化的必要条件：显示偏好的诞生}

设 $C_R(\cdot)$ 可被某偏好 $\succeq$ 生成（即理性化）。观察 $C_R(A)$：
\begin{itemize}
  \item $a\in C_R(A)$（$a$ 被选中了）$\iff a\succeq z$ 对所有 $z\in A$；
  \item $b\in A$（$b$ 当时也在场，本可被选择）。
\end{itemize}
于是可以\textbf{推断} $a\succeq b$——这就是“显示偏好”：行为把偏好“泄露”了出来。

再取另一个集合 $B\subseteq X$ 使 $a\in B$ 且 $b\in C_R(B)$
（$b$ 在 $B$ 中被选中 $\iff b\succeq z$ 对所有 $z\in B$），又可推断 $b\succeq a$。
两式合并：$a\sim b$，且由传递性 $a\in C_R(B)$。
\textbf{理性选择对跨集合行为施加了可检验的约束}——这正是下一节公理化的出发点。

\begin{example}[理性化判别表（务必手算一遍）]\label{ex:judgetable}
只知道 $C_R$ 在两个集合上的值：$A=\{a,b\}$ 与 $B=\{a,b,c\}$。
逐行判断该观察\textbf{是否可能被理性化}，并给出理由：
\begin{center}
\small
\begin{tabular}{cccp{5.6cm}}
\toprule
$C_R(\{a,b\})$ & $C_R(\{a,b,c\})$ & 可理性化？ & 推断 / 排除理由 \\
\midrule
$\{a\}$ & $\{c\}$ & $\checkmark$ & $c\succ a\succ b$（$c$ 全场击败 $a$，$a$ 击败 $b$） \\
$\{a\}$ & $\{a\}$ & $\checkmark$ & $a\succ b$、$a\succ c$；$b$ 与 $c$ 的关系未知 \\
$\{a,b\}$ & $\{c\}$ & $\checkmark$ & $c\succ a\sim b$ \\
$\{c\}$ & $\{c\}$ & $\times$ & $c\notin\{a,b\}$，违反选择规则的基本要求 $C_R(A)\subseteq A$ \\
$\varnothing$ & $\{c\}$ & $\times$ & 理性化要求选择非空（定理 \ref{thm:nonempty}），$a$ 与 $b$ 无法比较出最优 \\
$\{b\}$ & $\{a\}$ & $\times$ & 由 $\{a,b\}$ 得 $b\succ a$；由 $\{a,b,c\}$ 得 $a\succeq b$——反转 \\
$\{a\}$ & $\{a,b\}$ & $\times$ & 由 $\{a,b\}$ 得 $a\succ b$；由 $\{a,b,c\}$ 得 $a\sim b$——反转 \\
\bottomrule
\end{tabular}
\end{center}
读法：前三行是\textbf{相容}的——每条观察都能翻译成一条偏好推断，
且推断之间不打架；后四行各死于一种\textbf{结构性矛盾}：
违反 $C_R(A)\subseteq A$（第 4 行）、违反非空性（第 5 行）、
跨集合反转（第 6、7 行——这正是下一节 WARP/HARP 要禁止的情形）。
\textbf{方法}：把每条观察翻译成“$\succeq$ 记账”，再看账本是否自洽——
这个“记账”动作形式化后就是下面的 $R_C$。
\end{example}

形式化地：

\begin{definition}[直接与间接显示偏好关系]\label{def:RC}
对 $\mathcal{K}$ 上的选择函数 $C$，定义\textbf{直接显示偏好关系} $R_C$ 与
\textbf{直接显示严格偏好关系} $R_C^S$：
\begin{align}
R_C &= \{(x,y)\in X\times X \mid \exists\, S\in\mathcal{K}:\ x\in C(S)\ \&\ y\in S\},\\
R_C^S &= \{(x,y)\in X\times X \mid \exists\, S\in\mathcal{K}:\ x\in C(S)\ \&\ y\in S\setminus C(S)\}.
\end{align}
\textbf{间接显示偏好关系}为 $R_C$ 的传递闭包 $T(R_C)$；间接显示严格偏好为 $T(R_C^S)$。
\end{definition}

\begin{example}[手算算例：三步算出 $R_C$、$R_C^S$、$T(R_C)$ 并检验公理]\label{ex:workedRC}
设 $X=\{x,y,z\}$，定义域 $\mathcal{K}=\{S_1,S_2,S_3\}$，其中
$S_1=\{x,y\}$，$S_2=\{y,z\}$，$S_3=\{x,y,z\}$；观察到的选择为
\[
C(S_1)=\{x\},\qquad C(S_2)=\{y\},\qquad C(S_3)=\{z\}.
\]

\textbf{第一步：列 $R_C$（被选者对在场者逐一记账）。}
\begin{align*}
S_1:&\ x\in C(S_1),\ x,y\in S_1 \ \Rightarrow\ (x,x),(x,y)\in R_C;\\
S_2:&\ y\in C(S_2),\ y,z\in S_2 \ \Rightarrow\ (y,y),(y,z)\in R_C;\\
S_3:&\ z\in C(S_3),\ x,y,z\in S_3 \ \Rightarrow\ (z,x),(z,y),(z,z)\in R_C.
\end{align*}
故 $R_C=\{(x,x),(x,y),(y,y),(y,z),(z,x),(z,y),(z,z)\}$。

\textbf{第二步：列 $R_C^S$（只记“在场却落选”的）。}
$S_1$ 中 $y$ 落选 $\Rightarrow (x,y)\in R_C^S$；
$S_2$ 中 $z$ 落选 $\Rightarrow (y,z)\in R_C^S$；
$S_3$ 中 $x,y$ 落选 $\Rightarrow (z,x),(z,y)\in R_C^S$。
故 $R_C^S=\{(x,y),(y,z),(z,x),(z,y)\}$（注意：三个元素\textbf{两两}都被显示
严格偏好过——这正是选择过于“决绝”的后果）。

\textbf{第三步：取闭包、检验公理。}
$T(R_C)$：由 $x\to y\to z$、$y\to z\to x$、$z\to x\to y$ 等链，
$X\times X$ 全部九对都在 $T(R_C)$ 中。
\begin{itemize}
  \item \textbf{WARP}：$(x,y)\in R_C^S$ 要求 $(y,x)\notin R_C$——检查 $R_C$，
  无 $(y,x)$、$(z,y)$、$(x,z)$、$(y,z)$ 反向……逐对核验：$(x,y)\in R_C^S$
  而 $(y,x)\notin R_C$ $\checkmark$；$(y,z)\in R_C^S$ 而 $(z,y)\in R_C$ $\times$！
  $z$ 在 $S_2$ 中落选于 $y$（$(y,z)\in R_C^S$），却又在 $S_3$ 中于 $y$ 在场时
  被选（$(z,y)\in R_C$）——\textbf{WARP 已破}。
\end{itemize}
换一组数据体会“WARP 成立而 SARP 失败”：改 $C(S_3)=\{x,z\}$，其余不变。
重算得 $R_C^S=\{(x,y),(y,z)\}$（$S_3$ 中 $x,z$ 同被选，只产生对 $y$ 的
严格显示偏好），$R_C=R_C^S\cup\{(x,x),(y,y),(z,z),(x,z),(z,x)\}$。
$T(R_C^S)\ni (x,z)$（经 $y$），而 $(z,x)\in R_C$——故 \textbf{SARP 失败}
（$(x,z)\in T(R_C^S)$ 但 $(z,x)\in R_C$）。WARP 呢？逐对检查 $R_C^S$ 的两对：
$(x,y)$ 的反向 $(y,x)\notin R_C$ $\checkmark$；$(y,z)$ 的反向 $(z,y)\notin R_C$ $\checkmark$，
\textbf{WARP 成立}。结论：WARP 只看“一步反转”，管不住经中间项的
“两步绕行”——这正是要把 $R_C^S$ 升级为 $T(R_C^S)$ 的理由。
\end{example}

% ============================================================
\section{WARP、SARP、HARP 三大公理}
% ============================================================

\subsection{公理的陈述（Suzumura 形式）}

\begin{definition}[显示偏好三公理]\label{def:warp}
$\mathcal{K}$ 上的选择函数 $C$ 满足：
\begin{description}
  \item[WARP（弱显示偏好公理）] $\forall x,y\in X:\ (x,y)\in R_C^S \Rightarrow (y,x)\notin R_C$。
  即：若 $x$ 曾被显示严格优于 $y$（$x$ 被选而 $y$ 在场落选），则绝不能再观察到 $y$ 在 $x$ 在场时被选。
  \item[SARP（强显示偏好公理）] $\forall x,y\in X:\ (x,y)\in T(R_C^S) \Rightarrow (y,x)\notin R_C$。
  即把 WARP 的“直接”严格显示偏好升级为\textbf{链式}（间接）严格显示偏好。
  \item[HARP（Houthakker 显示偏好公理）] $\forall x,y\in X:\ (x,y)\in T(R_C) \Rightarrow (y,x)\notin R_C^S$。
  即：若 $x$ 经弱显示偏好链间接不弱于 $y$，则 $y$ 绝不能被显示严格优于 $x$。
\end{description}
\end{definition}

\begin{figure}[H]
\centering
\begin{tikzpicture}[scale=0.85,>=Stealth]
  \draw[axisline] (0,0) -- (7,0) node[below right] {$x_1$};
  \draw[axisline] (0,0) -- (0,5.5) node[above left] {$x_2$};
  \fill[fillsoft] (0,0) -- (6,0) -- (0,4) -- cycle;
  \draw[very thick,AxisBlue] (6,0) -- (0,4) node[above right,sloped,pos=0.55] {\small $B(p,m)$};
  \fill[nodeblue] (4,1.33) circle (2.2pt) node[below right] {$x$};
  \fill[black] (2,2.67) circle (1.8pt) node[above left] {$y$};
  \draw[thick,dashed,DeepBlue] (5.5,0) -- (0,5) node[above right,sloped,pos=0.6] {\small $B(p',m')$};
  \node[align=left,font=\scriptsize] at (1.6,0.7) {观察1：$x$ 被选，\\$y$ 可行};
  \node[align=left,font=\scriptsize] at (4.6,3.5) {观察2：$y$ 被选，\\$x$ 可行};
  \draw[arrblue] (3.6,1.5) -- (2.9,2.2) node[midway,below right,font=\scriptsize] {违反 WARP};
\end{tikzpicture}
\caption{WARP 违反：$x$ 曾显示优于 $y$，后来又观察到 $y$ 在 $x$ 可行时被选}
\label{fig:warp_violation}
\end{figure}
\figexplan{浅灰三角为预算集 $B(p,m)$，实斜边为其预算线；虚线为另一期预算线 $B(p',m')$。观察 1 中 $x$ 被选中而 $y$ 可行，故 $x$ 显示优于 $y$；观察 2 中 $y$ 被选中而 $x$ 可行，故 $y$ 又显示优于 $x$。箭头标出的两次反转即构成对 WARP 的违反。}

\begin{example}[家庭支出数据的 GARP 检验：一个实证案例]
Varian（1982）将显示偏好理论转化为可计算的检验：给定一个家户在 $T$ 期的消费数据 $\{(p^t,x^t)\}_{t=1}^T$，构造直接显示偏好 $x^s\succeq^R x^t \iff p^s\cdot x^s\ge p^s\cdot x^t$，再检查是否存在违反 GARP 的链式反转。Blundell、Browning 与 Crawford（2003, 2008）对英国家庭支出调查（FES）的应用显示：在食品、休闲、交通等大类上，绝大多数家户的数据满足 GARP；但当消费类别划分过细或包含耐用品时，违反率显著上升——这提示模型设定（商品聚合、时间跨度）本身就是可被证伪的实证命题。
\end{example}

\begin{proposition}[逻辑强度]
\[
(\mathrm{HARP}) \Rightarrow (\mathrm{SARP}) \Rightarrow (\mathrm{WARP}),
\]
且反向均不成立。
\end{proposition}

\begin{proof}[证明要点]
由定义 $R_C^S\subseteq R_C$，两边取传递闭包得 $T(R_C^S)\subseteq T(R_C)$，
故 HARP 的前提更“宽”而结论相同，故 HARP $\Rightarrow$ SARP；
SARP $\Rightarrow$ WARP 因 $R_C^S\subseteq T(R_C^S)$。反向失败的反例见 §2.6.3。
\end{proof}

\begin{remark}[历史坐标]
Samuelson (1938) 在竞争预算域 $\mathcal{K}_{CB}$ 上对单值需求函数提出 WARP；
Houthakker (1950, \emph{Economica}) 提出\textbf{半传递性公理}（semitransitivity），
即预算域上的 SARP 雏形，并首次证明：满足该公理时存在效用函数 $u$ 使选择可由
$u$ 在 $B(\boldsymbol{p},M)$ 上最大化来理性化——这是“选择即理性化”理论的第一个重大成果。
Houthakker 本人回忆，Samuelson 收到他的论文时说：“他不仅补上了我缺失的证明，
而且独立提出了正是我寄望于挽局的那个强公理。”
随后分两路推广：Arrow (1959) 与 Sen (1971) 走\textbf{有限域}多值选择方向；
Richter (1966, 1971) 与 Hansson (1968) 走\textbf{任意域}方向。
\end{remark}

\begin{remark}[Houthakker 半传递性公理的原文形式（预算域版）]\label{rem:houthakker}
在竞争预算域 $\mathcal{K}_{CB}$ 上，记第 $k$ 期价格—收入为 $(\boldsymbol{p}^k,M^k)$，
预算集 $B(\boldsymbol{p}^k,M^k)=\{x\in\R^n_+\mid \boldsymbol{p}^k\cdot x\le M^k\}$，
需求 $x^k=c(B(\boldsymbol{p}^k,M^k))$（单值需求函数，$c$ 满足预算平衡
$\boldsymbol{p}^k\cdot x^k=M^k$）。Houthakker (1950) 的\textbf{半传递性公理}：
对任意有限观察序列 $k=0,1,\dots,K$，
\[
\boldsymbol{p}^{k-1}\cdot x^{k-1}\ge \boldsymbol{p}^{k-1}\cdot x^k
\ \ (k=1,\dots,K)\ \Longrightarrow\
\boldsymbol{p}^{K}\cdot x^{K}<\boldsymbol{p}^{K}\cdot x^{0},
\]
即：若 $x^1$ 在 $x^0$ 的价格下买得起（$x^0$ 直接显示不弱于 $x^1$），$x^2$ 在
$x^1$ 的价格下买得起……沿链 $x^0\succeq^D x^1\succeq^D\cdots\succeq^D x^K$，
则在第 $K$ 期 $x^0$ 必须\textbf{买不起}（$x^K$ 间接显示严格优于 $x^0$）。
对比 Samuelson (1938) 的 WARP 只禁止\textbf{两步}反转（$K=1$ 情形），
半传递性把禁令推广到\textbf{任意有限链}——这就是预算域上的 SARP。
Houthakker 证明：需求函数满足半传递性 $\iff$ 存在效用函数 $u$ 使
$c(B(\boldsymbol{p},M))=\arg\max\{u(x)\mid \boldsymbol{p}\cdot x\le M\}$。
其证明构造性地把显示偏好关系扩张为序（用了 Szpilrajn 式论证），
再对扩张出的序调用效用表示定理——本节定理 \ref{thm:szpilrajn} 正是在这一步登场的。
\end{remark}

\subsection{成对表述、有限数据与 \texorpdfstring{$\widehat{C}_R$}{Ĉ\_R}}

显示偏好公理族还有一种更紧凑的\textbf{成对表述}，并便于直接处理“数据有限”的现实约束。

\begin{definition}[HARP 与 WARP（成对形式）]\label{def:slideswarp}
\begin{description}
  \item[HARP（成对形式）] 显示偏好选择规则 $C_R: 2^X\to 2^X$ 满足 HARP，
  当且仅当：对任意 $a,b\in X$ 与任意 $A,B\subseteq X$，若
  \[
  \{a,b\}\subseteq A\ \text{且}\ a\in C_R(A);\qquad
  \{a,b\}\subseteq B\ \text{且}\ b\in C_R(B),
  \]
  则 $a\in C_R(B)$（并由对称性 $b\in C_R(A)$）。\\
  读法：$a$ 在 $b$ 在场时赢过、$b$ 在 $a$ 在场时也赢过——那么二人必须
  \textbf{共享每一次胜利}（否则如 §2.3.4 的推断，$a\sim b$ 被破坏）。
  \item[WARP（成对形式）] 同一条件，但只对定义域 $\mathcal{B}\subseteq 2^X$
  内的 $A,B\in\mathcal{B}$ 要求。\\
  即口诀：\textbf{HARP $=$ WARP $+$ $\mathcal{B}=2^X$}
  （全域上的 WARP 就是 HARP）。
\end{description}
\end{definition}

\begin{remark}[与 Suzumura 形式的对应]
成对形式排除的正是“$a$ 赢过 $b$、$b$ 又赢过 $a$ 却互不共享”的两步反转，
与定义 \ref{def:warp} 的 WARP（$R_C^S$ 无反转）逐点对应：
$a\in C_R(A)$、$b\in A\setminus C_R(A)$ 即 $(a,b)\in R_C^S$。
在\textbf{全域} $2^X$ 上，成对形式还蕴含链式形式——任何显示偏好链都可
用二元、三元集逐步接通（§2.5 定理 \ref{thm:harp} 的传递性证明正是这么做的），
故成对形式的 HARP 与链式形式的 HARP 等价；在\textbf{非全域}上，
成对 WARP 严格弱于 Suzumura 的链式公理——见下例。
\end{remark}

\begin{remark}[有限数据的两种缺口]
现实的观察从来不完备，缺口有两种：
\begin{itemize}
  \item \textbf{只看到一个被选项}：观察值是某个 $\widetilde{C}_R(A)\in C_R(A)$
  ——选择集里还有哪些并列最优，数据不告诉我们；
  \item \textbf{只看到部分机会集}：观察值是 $\widehat{C}_R:\mathcal{B}\to 2^X$，
  其中 $\mathcal{B}\subsetneq 2^X$，在 $\mathcal{B}$ 上 $\widehat{C}_R(A)=C_R(A)$，
  域外一无所知。
\end{itemize}
WARP、GARP 等公理正是为 $\widehat{C}_R$ 这类\textbf{部分定义}的选择规则设计的；
GARP（广义显示偏好公理）是有限数据下理性化的充要条件——
其预算域版本即 Afriat 定理（§2.7.5，式 \eqref{eq:afriat}）。
\end{remark}

\begin{example}[WARP 不足以理性化（对照例 \ref{ex:WARPnotAR}）]\label{ex:warpnotsuff}
设 $X=\{a,b,c\}$，定义域仅含三个二元集
$\mathcal{B}=\bigl\{\{a,b\},\{b,c\},\{c,a\}\bigr\}$，观察为
\[
\widehat{C}_R(\{a,b\})=\{a\},\qquad
\widehat{C}_R(\{b,c\})=\{b\},\qquad
\widehat{C}_R(\{c,a\})=\{c\}.
\]
\textbf{逐对检验 WARP}：每个集合里落选者从此\textbf{不再出现于任何其他集合的选择集中}
——$b$ 在 $\{a,b\}$ 落选，但 $b$ 只在 $\{b,c\}$ 中被选而那里 $a$ 不在场……
严格地：成对反转要求存在 $A,B\in\mathcal{B}$ 使 $\{x,y\}\subseteq A\cap B$，
而这里任意两集合只交一个元素，反转\textbf{无从发生}，WARP 平凡成立。

\textbf{但不可理性化}：三条观察翻译成账本 $a\succ b$、$b\succ c$、$c\succ a$——
严格偏好环。也可换用 HARP 语言：把 $\widehat{C}_R$ 想成某个全域
$C_R: 2^X\to 2^X$ 的限制，则无论 $C_R(\{a,b,c\})$ 取什么非空值，
都会与某条二元观察构成反转（取 $\{a\}$ 则与 $\{c,a\}$ 的观察冲突，余类似）——
\textbf{不存在与 HARP 相容的 $C_R(\{a,b,c\})$}。
这正是“WARP（成对）$\not\Rightarrow$ SARP（链式）”的最小实例：
链 $a\succ b\succ c\succ a$ 有三节，成对检查每节都合法。
\end{example}

\begin{figure}[H]
\centering
\begin{tikzpicture}[scale=1.1,>=Stealth]
  \node[nodeblue] (a) at (0,0) {$a$};
  \node[nodeblue] (b) at (3,0) {$b$};
  \node[nodeblue] (c) at (1.5,2.6) {$c$};
  \draw[->,thick,AxisBlue] (a) -- (b) node[midway,below] {$a\succ b$};
  \draw[->,thick,AxisBlue] (b) -- (c) node[midway,right] {$b\succ c$};
  \draw[->,thick,AxisBlue] (c) -- (a) node[midway,left] {$c\succ a$};
\end{tikzpicture}
\caption{三元严格偏好环：每对二元比较都合法，但整体不可被理性化}
\label{fig:cycle}
\end{figure}
\figexplan{三个节点与三条箭头构成 $a\succ b$、$b\succ c$、$c\succ a$ 的循环。每一对二元比较本身都合法，但三者合起来不存在最大元，故整体无法被任何完备传递的偏好关系（因而也无法被效用函数）理性化——这说明逐对检查不足以验证理性。}

\subsection{Hansson 形式：C-连通表述}

Suzumura（Definition 1.6）给出等价但更便于构造反例的表述。称 $(S_1,S_2)$ 为
$\mathcal{K}$ 中的 \textbf{$C$-连通对}，若 $S_1\cap C(S_2)\neq\varnothing$ 且
$C(S_1)\cap S_2\neq\varnothing$；称 $(S_1,\dots,S_n)$ 为 \textbf{$C$-连通序列}，若
$S_k\cap C(S_{k+1})\neq\varnothing$（$k=1,\dots,n-1$）且 $S_n\cap C(S_1)\neq\varnothing$。

\begin{itemize}
  \item \textbf{WARP}：对任意 $C$-连通对，$S_1\cap C(S_2)=C(S_1)\cap S_2$；
  \item \textbf{SARP}：对任意 $C$-连通序列，存在某个 $k$ 使 $S_k\cap C(S_{k+1})=C(S_k)\cap S_{k+1}$；
  \item \textbf{HARP}：对任意 $C$-连通序列，\textbf{所有} $k$ 都满足 $S_k\cap C(S_{k+1})=C(S_k)\cap S_{k+1}$。
\end{itemize}

“连通对的两侧交集相等”等价于“无反转”：若 $x\in S_1\cap C(S_2)$ 而 $y\in C(S_1)\cap S_2$ 且 $x\neq y$，
则 $x$ 在 $S_2$ 中击败 $y$ 又在 $S_1$ 中被 $y$ 击败——恰是 WARP 禁止的反转。
两种表述的等价性在 Suzumura 书 Appendix 1.A 中逐步证明（关键引理：
WARP $\iff$ (1.1) 式、SARP $\iff$ (1.2) 式）。

\subsection{Richter 的一致性公理}

除上述三公理外，Richter (1966; 1971) 还提出两条以 $R_C$ 陈述的一致性条件：
\begin{description}
  \item[DRC（直接显示一致性）] $\forall S\in\mathcal{K},\ \forall x\in S:\
  [\forall y\in S:\ (x,y)\in R_C \Rightarrow x\in C(S)]$；
  \item[TCC（传递闭包一致性）] $\forall S\in\mathcal{K},\ \forall x\in S:\
  [\forall y\in S:\ (x,y)\in T(R_C) \Rightarrow x\in C(S)]$。
\end{description}
读法：若 $x$ 对 $S$ 中每个 $y$ 都（直接或间接地）显示不弱，则 $x$ 必须真的被选中。
它们是 G-理性化的候选刻画条件（见 §2.6.2 的 G-normality）。

\subsection{Sen 的条件 \texorpdfstring{$\alpha$、$\beta$、$\gamma$}{α、β、γ}：选择一致性的另一语言}

与显示偏好三公理平行，Sen (1971, \emph{Econometrica}) 在《Choice Functions and
Revealed Preference》中以\textbf{集合包含语言}陈述了一组选择一致性条件
（记号沿用 Sen 原文）：

\begin{definition}[Sen 条件 $\alpha$、$\beta$、$\gamma$]\label{def:senabc}
设 $C$ 为 $\mathcal{K}$ 上的选择函数，$S_1,S_2\in\mathcal{K}$。
\begin{description}
  \item[条件 $\alpha$（收缩一致性）] 若 $S_1\subseteq S_2$，则
  $C(S_2)\cap S_1\subseteq C(S_1)$。\\
  读法：大菜单里的最优项，进了小菜单若还在场，仍须是最优
  （“世界冠军必是全国冠军”）。
  \item[条件 $\beta$（扩张一致性）] 若 $S_1\subseteq S_2$ 且
  $x,y\in C(S_1)$，则 $x\in C(S_2)\iff y\in C(S_2)$。\\
  读法：小菜单里并列最优的两者，到大菜单里仍须同进退
  （不允许“竞争加剧反而拆散并列”）。
  \item[条件 $\gamma$（并集一致性）] 对任意族 $\{S_j\}\subseteq\mathcal{K}$，
  若 $x\in C(S_j)$ 对每个 $j$ 成立，且 $\bigcup_j S_j\in\mathcal{K}$，则
  $x\in C(\bigcup_j S_j)$。\\
  读法：在每个分赛场都夺冠，合并后的总决赛也应夺冠。
\end{description}
\end{definition}

\begin{theorem}[Sen 1971 的核心结果]
在包含所有有限子集的定义域上：
\begin{enumerate}
  \item $C$ 可被\textbf{完备}二元关系理性化（G-理性化）$\iff$ 条件 $\alpha$ + 条件 $\gamma$；
  \item $C$ 可被\textbf{序}理性化（完全理性）$\iff$ 条件 $\alpha$ + 条件 $\beta$
  （此时 $\gamma$ 自动成立）。
\end{enumerate}
\end{theorem}

\begin{remark}[$\alpha/\beta/\gamma$ 与 WARP 家族的关系]
两套语言在标准域上互通：$\alpha+\beta\Rightarrow$ WARP；而 WARP + 非空值 +
有限域闭合性 $\Rightarrow \alpha+\beta$。但\textbf{在任意域上它们分道扬镳}——
$\alpha$、$\beta$、$\gamma$ 的合取既不能蕴含、也不被蕴含于 Suzumura 的
同余公理体系（这正是 §2.6 需要 WCA/SCA 的原因：同余语言直接对显示偏好关系记账，
不受域结构限制）。Sen 的表述胜在直观、易向非专业听众解释（社会选择文献多用之）；
Suzumura--Richter 的表述胜在一般性。学有余力者应\textbf{双语互译}：
把每个 $\alpha/\beta$ 反例改写成 $R_C$ 记账语言，再反过来。
\end{remark}

% ============================================================
\section{枢纽定理：HARP 与理性化等价}
% ============================================================

这是本章第一座高峰。

\begin{theorem}[HARP 是理性化的充分必要条件]\label{thm:harp}
设显示偏好选择规则 $C_R: 2^X\to 2^X$ 非空值（除 $C_R(\varnothing)$ 外）。
则 $C_R(\cdot)$ 满足 HARP，当且仅当存在理性偏好关系 $\succeq$（完备且传递）使
$C_R(\cdot)=C(\cdot,\succeq)$。
\end{theorem}

\begin{proof}
\textbf{（$\Rightarrow$ 必要性）}有理化 $\Rightarrow$ HARP：如前节“必要条件”论证，
有理化的选择规则必然不产生显示偏好反转，即 HARP 成立。

\textbf{（$\Leftarrow$ 充分性）}设 $C_R(\cdot)$ 满足 HARP。构造\textbf{显示偏好关系} $\succeq_C$：
对任意 $x,y$，令
\[
x\succeq_C y \iff \text{存在某个}\ A\subseteq X\ \text{使}\ y\in A\ \text{且}\ x\in C_R(A).
\]
需证 $\succeq_C$ 完备、传递，且生成 $C_R$（即 $C_R(\cdot)=C(\cdot,\succeq)$）。

\textbf{第一步：完备性。}
$C_R$ 非空值，故 $C_R(\{x,y\})\neq\varnothing$：要么 $x\in C_R(\{x,y\})$，要么
$y\in C_R(\{x,y\})$。于是 $x\succeq_C y$ 或 $y\succeq_C x$。

\textbf{第二步：传递性。}
设 $x\succeq_C y$ 且 $y\succeq_C z$。考察三元素集 $C_R(\{x,y,z\})\neq\varnothing$，
以下三种情形必居其一：
\begin{enumerate}
  \item $x\in C_R(\{x,y,z\})$：直接得 $x\succeq_C z$；
  \item $y\in C_R(\{x,y,z\})$：由 $x\succeq_C y$ 与 HARP（$y$ 在 $x$ 可得的集合中被选中
  而 $x$ 在场，若 $x$ 落选则反转），必有 $x\in C_R(\{x,y,z\})$，归约到情形 1，故 $x\succeq_C z$；
  \item $z\in C_R(\{x,y,z\})$：由 $y\succeq_C z$ 与 HARP 得 $y\in C_R(\{x,y,z\})$，
  归约到情形 2，从而 $x\succeq_C z$。
\end{enumerate}
三种情形均得 $x\succeq_C z$，传递性成立。

\textbf{第三步：生成性。}需证对任意 $B\subseteq X$：
$x\in C_R(B) \iff x\succeq_C y$ 对所有 $y\in B$。
\begin{itemize}
  \item（$\Rightarrow$）$x\in C_R(B)$，由 $\succeq_C$ 的构造即 $x\succeq_C y$ 对所有 $y\in B$；
  \item（$\Leftarrow$）设 $x\succeq_C y$ 对所有 $y\in B$。由非空值性存在某
  $y^*\in C_R(B)$；由假设 $x\succeq_C y^*$，再由 HARP（$y^*$ 被选而 $x$ 在场，
  $x$ 不能落选），得 $x\in C_R(B)$。
\end{itemize}
于是 $C_R(\cdot)=C(\cdot,\succeq)$，且 $\succeq_C$ 完备传递，证毕。
\end{proof}

\begin{remark}[这个定理为什么重要]
它给出选择理论的一个\textbf{完全可检验}的刻画：要判断一批观察到的选择是否“理性”，
不需要猜偏好是什么，只需检验选择数据本身是否满足 HARP。
\textbf{理性（不可观察）$\iff$ HARP（纯数据性质）}——这是显示偏好方法论的巅峰成果，
也是 Richter (1966)、Hansson (1968)、Suzumura (1977) 的主定理。
\end{remark}

\begin{remark}[教学提示：HARP 的图形]
用两个相交的机会集 $B, B'$ 演示 HARP 的违反：
$x\in C(B)\cap B'$（$x$ 在 $B$ 中被选且在 $B'$ 中可得），$x'\in C(B')\cap B$
（$x'$ 在 $B'$ 中被选且在 $B$ 中可得），$x\neq x'$。
此时 $x$ 显示不弱于 $x'$ 且 $x'$ 显示不弱于 $x$——若理性化存在，二者应无差异并同时出现在
$C(B)$ 与 $C(B')$ 中；观察到的却是互斥选择，即 HARP（及其推论 $B\cap C(B')=C(B)\cap B'$）被破坏。
课堂自测：若 $C_R(\{a,b\})=\{b\}$，能对 $C_R(\{a,b,c\})$ 与 $C_R(\{a,b\})$ 的关系下什么结论？
\end{remark}

% ============================================================
\section{任意定义域上的一般理论：同余公理与理性等级}
% ============================================================

Suzumura 的第二章（Chapter 1, Essay 1）把上述理论推广到\textbf{任意定义域}
$\mathcal{K}_A$——不要求包含所有有限集或预算集。“冷风穿过无结构的集合”（Nigel Howard 语），
一般域上定理更少、反例更多，这正是一般理论的价值所在。

\subsection{同余公理：WCA 与 SCA}

\begin{definition}[Richter--Sen 同余公理]\label{def:congruence}
$\mathcal{K}$ 上的选择函数 $C$ 满足：
\begin{itemize}
  \item[(a)] \textbf{弱同余公理（WCA）}：对任意 $S\in\mathcal{K}$，
  $(x,y)\in R^*$，$x\in S$ 且 $y\in C(S)$ 蕴含 $x\in C(S)$；
  \item[(b)] \textbf{强同余公理（SCA）}：对任意 $S\in\mathcal{K}$，
  $(x,y)\in T(R^*)$，$x\in S$ 且 $y\in C(S)$ 蕴含 $x\in C(S)$。
\end{itemize}
（此处 $R^*$ 即 $C$ 的直接显示偏好关系 $R_C$。）
\end{definition}

读法：$y$ 被选中、$x$（直接或间接地）显示不弱于 $y$、$x$ 又在场——那么 $x$ 也必须被选中。
同余公理把“显示不弱”封闭到选择集内部。

\begin{theorem}[公理等价]\label{thm:equiv}
在任意定义域上：
\[
(\mathrm{WARP}) \iff (\mathrm{WCA}), \qquad (\mathrm{HARP}) \iff (\mathrm{SCA}).
\]
\end{theorem}

\begin{proof}[证明骨架（Suzumura, Theorem 1.1）]
\textbf{(i) WARP $\iff$ WCA}，分三步：
\begin{enumerate}
  \item $\mathrm{WARP} \Rightarrow (R^*\subseteq R_*)$：若不然，存在 $(x,y)\in R^*\setminus R_*$，
  即有 $S,S'$ 使 $x\in C(S), y\in S, x\in S', x\notin C(S'), y\in C(S')$——
  于是 $x\in S'\cap C(S)$、$y\in C(S')\cap S$、$x\notin C(S')\cap S$，恰违反 WARP。
  \item $(R^*\subseteq R_*) \Rightarrow \mathrm{WCA}$：$(x,y)\in R^*$、$x\in S$、$y\in C(S)$
  推出 $(x,y)\in R_*$，由 $R_*$ 的定义即 $x\in C(S)$。
  \item $\mathrm{WCA} \Rightarrow \mathrm{WARP}$：对 $C$-连通对 $(S_1,S_2)$，
  取 $x\in S_1\cap C(S_2)$、$y\in C(S_1)\cap S_2$。由 $x\in C(S_2)$、$y\in S_2$ 得
  $(x,y)\in R^*$；又 $x\in S_1$、$y\in C(S_1)$，由 WCA 得 $x\in C(S_1)$。
  故 $S_1\cap C(S_2)\subseteq C(S_1)\cap S_2$，对称论证给出反向包含，等式成立。
\end{enumerate}
\textbf{(ii) HARP $\iff$ SCA}，分两步：
\begin{enumerate}
  \item $\mathrm{HARP} \Rightarrow \mathrm{SCA}$：设 $(x,y)\in T(R^*)$、$x\in S$、$y\in C(S)$。
  若 $(x,y)\in R^*$，因 HARP $\Rightarrow$ WARP $\iff$ WCA 即得 $x\in C(S)$；
  否则存在链条 $\{z^k\}$、$\{S_k\}$ 把 $x$ 与 $y$ 接通，沿链反复用 HARP
  （得 $S_k\cap C(S_{k+1})=C(S_k)\cap S_{k+1}$）把“选中”一步步传递到 $x$。
  \item $\mathrm{SCA} \Rightarrow \mathrm{HARP}$：对 $C$-连通序列取 $z^k\in S_k\cap C(S_{k+1})$、
  $z^n\in S_n\cap C(S_1)$，用 SCA 逐段证明 $S_1\cap C(S_2)=C(S_1)\cap S_2$ 及一般段等式。
\end{enumerate}
\end{proof}

\begin{remark}
SCA $\Rightarrow$ WCA 显然（$R^*\subseteq T(R^*)$），反向不成立——见例 \ref{ex:WCAnotSCA}。
\end{remark}

\subsection{三大刻画定理与 G/M-正规性}

\begin{theorem}[Arrow--Sen 结果（$\alpha$）]
若定义域 $\mathcal{S}$ 包含 $X$ 的所有二元子集与三元子集，则
SCA、WCA、SARP、WARP \textbf{互为等价的}完全理性（full rationality）的充分必要条件。
\end{theorem}

\begin{theorem}[Richter--Hansson 结果（$\beta$）]
不对定义域内容作任何限制：SCA 与 HARP 互为等价，且是完全理性的充分必要条件。
\end{theorem}

\begin{theorem}[Richter 结果（$\gamma$）]
G-正规性（G-normality）是 G-理性化的充分必要条件。
\end{theorem}

其中正规性概念为：记 $G^*(S)=G(S,R^*)$、$M^*(S)=M(S,R^*)$ 为按显示偏好关系计算的
最优点集与最大点集。称 $C$ 是 \textbf{G-正规}（\textbf{M-正规}）的，若 $C(S)=G^*(S)$
（$C(S)=M^*(S)$）对所有 $S$ 成立。一个基本链条是
\[
C(S)\subseteq G^*(S)\subseteq M^*(S) \quad\text{对所有}\ S,
\]
左包含因 $x\in C(S)$ 时对一切 $y\in S$ 有 $(x,y)\in R^*$；右包含即注 \ref{rem:GM}。
于是（Theorem 1.3）WARP $\Rightarrow$ $M^*(S)\subseteq C(S)$，即在 WARP 下选择函数恰为
M-正规，从而为 M-理性。

\begin{remark}[三大结果的分工]
（$\alpha$）代价是定义域假设太强（要包含全部二元、三元组），限制了普适性；
（$\beta$）对定义域\textbf{不作限制}，是一般理论的主定理；
（$\gamma$）处理不要求完备性的 G-理性化。Suzumura 的贡献在于把 SARP、WARP、WCA、
M-正规性在这一般框架下的角色彻底厘清（Figure 1.1，见本章图 \ref{fig:logical}）。
\end{remark}

以下补出上文引用却略去证明的三个关键结果，均为本章的\textbf{高难度精读段落}。

\begin{theorem}[Suzumura, Theorem 1.3：WARP $\Rightarrow$ M-正规]\label{thm:13}
任意定义域上，若 $C$ 满足 WARP，则 $M^*(S)\subseteq C(S)$ 对所有 $S\in\mathcal{K}$
成立，从而 $C(S)=M^*(S)$（$C$ 为 M-正规）。
\end{theorem}

\begin{proof}[完整证明]
已知基本链条 $C(S)\subseteq G^*(S)\subseteq M^*(S)$，只需证反向包含
$M^*(S)\subseteq C(S)$。

任取 $x\in M^*(S)$。由最大点集定义，$x\in S$ 且
\begin{equation}
\forall y\in S:\ (y,x)\notin P(R^*). \tag{$\ast$}
\end{equation}
由定理 \ref{thm:equiv} 证明第 (i) 步，WARP $\Rightarrow R^*\subseteq R_*$
（其中 $(a,b)\in R_*$ 定义为：不存在 $S'$ 使 $a\in S'\setminus C(S')$ 且
$b\in C(S')$——即“$a$ 从未在 $b$ 被选时在场落选”）。

取任意 $z\in C(S)$（$C$ 非空值）。因 $x,z\in S$ 且 $z\in C(S)$，有
$(z,x)\in R^*$。若 $(z,x)\in P(R^*)$，即 $(z,x)\in R^*$ 且 $(x,z)\notin R^*$：
由 $(x,z)\notin R^*$，对一切含 $x,z$ 的集合 $x$ 都未在其被选时出现于 $z$ 在场
的场合——特别地，在任何 $S'$ 中 $z\in C(S')$ 且 $x\in S'$ 的场合……
直接地：$(x,z)\notin R^*$ 意味着不存在 $S'\in\mathcal{K}$ 使 $x\in C(S')$ 且
$z\in S'$。但这尚不足以定位 $x$ 是否在 $C(S)$ 中，改用 $R^*\subseteq R_*$：
由 $(z,x)\in R^*\subseteq R_*$，$R_*$ 的定义给出
“$z$ 从未在 $x$ 被选时在场落选”。现在用反证：设 $x\notin C(S)$。
则 $z\in C(S)$、$x\in S\setminus C(S)$，这正说明“$x$ 在 $z$ 被选时在场落选”，
即 $(z,x)\in R^{*S}\subseteq P(R^*)$——与 $(\ast)$（取 $y=z$）矛盾。
故 $x\in C(S)$，$M^*(S)\subseteq C(S)$ 得证。

合并基本链条即 $C(S)=G^*(S)=M^*(S)$ 对所有 $S$ 成立。
\end{proof}

\begin{theorem}[M-正规 $\Rightarrow$ 完备 G-理性化（Theorem 1.5 方向）]\label{thm:15}
若 $C$ 是 M-正规的（$C(S)=M^*(S)$ 对所有 $S$），则 $C$ 可被\textbf{完备}
二元关系 M-理性化（因而是完备 G-理性化）。
\end{theorem}

\begin{proof}[证明]
取候选理性化关系为 $R^0\equiv R^*\cup\Delta_X\cup\{(x,y)\mid x,y \text{ 在 } R^*
\text{ 下不可比}\}$ 的\textbf{完备化版本}——标准构造（Suzumura, p.\,36）：
\[
R^0 = R^*\ \cup\ \{(y,x)\in X\times X \mid (x,y)\notin R^*\}.
\]
即：凡 $R^*$ 未给出 $x R^* y$ 的，一律补记 $y R^0 x$。$R^0$ \textbf{完备}显然
（任给 $x,y$，要么 $(x,y)\in R^*\subseteq R^0$，要么 $(y,x)\in R^0$ 由补记条款）。

需证 $C(S)=M(S,R^0)$ 对所有 $S$ 成立。
\begin{itemize}
  \item[$\subseteq$：] 设 $x\in C(S)=M^*(S)$。对任意 $y\in S$，由 $(\ast)$ 式
  $(y,x)\notin P(R^*)$，即 $(y,x)\in R^*\Rightarrow(x,y)\in R^*$。
  若 $(y,x)\in R^*$ 则 $(x,y)\in R^*\subseteq R^0$；若 $(y,x)\notin R^*$ 则由
  补记条款 $(x,y)\in R^0$。两情形均得 $(x,y)\in R^0$，故 $x\in G(S,R^0)$。
  再由 $R^0$ 完备，$G(S,R^0)=M(S,R^0)$（完备关系下最优点集与最大点集重合：
  $x$ 是最大点 $\iff$ 无 $y$ 使 $(y,x)\in P(R^0)$，而完备性使
  $(x,y)\notin R^0\iff(y,x)\in P(R^0)$）。于是 $x\in M(S,R^0)$。
  \item[$\supseteq$：] 设 $x\in M(S,R^0)=G(S,R^0)$（完备性），即对所有 $y\in S$
  有 $(x,y)\in R^0$。若 $x\notin M^*(S)=C(S)$，则存在 $y\in S$ 使
  $(y,x)\in P(R^*)$，即 $(y,x)\in R^*$ 且 $(x,y)\notin R^*$。
  但 $(x,y)\in R^0$ 且 $(x,y)\notin R^*$ 只能来自补记条款，
  补记条款要求 $(y,x)\notin R^*$——矛盾。故 $x\in M^*(S)=C(S)$。
\end{itemize}
两向包含给出 $C(S)=M(S,R^0)$：$R^0$ 是 $C$ 的完备 M-理性化。
\end{proof}

\begin{theorem}[SARP $\Rightarrow$ 无环完备理性（Theorem 1.6 方向）]\label{thm:16}
若 $C$ 满足 SARP，则存在\textbf{完备且无环}的二元关系 $R^0$ 使
$C(S)=G(S,R^0)$ 对所有 $S\in\mathcal{K}$ 成立。
\end{theorem}

\begin{proof}[证明（$R^0$ 的显式构造）]
沿用定理 \ref{thm:15} 的完备化 $R^0=R^*\cup\{(y,x)\mid(x,y)\notin R^*\}$。
由 SARP $\Rightarrow$ WARP（命题 4.1）与定理 \ref{thm:13}，$C$ 为 M-正规；
再由定理 \ref{thm:15} 的证明，$C(S)=M(S,R^0)=G(S,R^0)$（$R^0$ 完备）。
剩下的核心任务是证 $R^0$ \textbf{无环}——这是 SARP 超出 WARP 的增量所在。

设若不然，存在严格环 $x^1,x^2,\dots,x^K\in X$（$K\ge 2$）：
$(x^k,x^{k+1})\in P(R^0)$（$k=1,\dots,K-1$）且 $(x^K,x^1)\in P(R^0)$。

\textbf{断言：}$(x^k,x^{k+1})\in P(R^0)$ 蕴含 $(x^k,x^{k+1})\in T(R^*)$ 且
$(x^{k+1},x^k)\notin T(R^*)$。
事实上，$P(R^0)$ 的一对 $(a,b)$ 只有两种来源：(i) $(a,b)\in P(R^*)$；或
(ii) 补记条款，即 $(a,b)\in R^0\setminus R^*$ 且 $(b,a)\notin R^*$，
而严格性 $(b,a)\notin R^0$ 又要求 $(a,b)\notin R^*$ 时 $(b,a)\notin R^*$——
情形 (ii) 给出的是“$R^*$-不可比对”被定向。对情形 (ii)，$(a,b)\in P(R^0)$
要求 $(b,a)\notin R^0$，但补记条款对不可比对是\textbf{对称写入}的
（$(b,a)\notin R^*\Rightarrow(a,b)\in R^0$ 且 $(a,b)\notin R^*\Rightarrow(b,a)\in R^0$），
故严格对只能来自情形 (i) 的链条闭包……严格地说，定义
$\bar R = T(R^*)\cup\{(y,x)\mid (x,y)\notin T(R^*)\}$ 为对 $T(R^*)$ 的同一完备化，
则 $\bar R$ 完备，且 $(a,b)\in P(\bar R)$ 当且仅当 $(a,b)\in T(R^*)$ 而
$(b,a)\notin T(R^*)$（补记条款在 $T(R^*)$ 下自反不成严格）。
于是上述严格环给出 $(x^k,x^{k+1})\in T(R^*)$ 全链且 $(x^1,x^K)\notin T(R^*)$。
但 $T(R^*)$ 传递：全链合成得 $(x^1,x^K)\in T(R^*)$——矛盾。

因此 $P(\bar R)$ 无环，即 $\bar R$ 完备无环且（与定理 \ref{thm:15} 同样的
双向包含论证，把 $R^*$ 换成 $T(R^*)$、并注意 SARP 正是“$T(R^{*S})$ 无反转”）
$C(S)=G(S,\bar R)$ 对所有 $S$ 成立。
\end{proof}

\begin{remark}[三个证明的方法论读法]
三个结果串成一条链：WARP 把选择\textbf{钉死}在显示偏好最大点上（Thm 1.3），
M-正规保证存在\textbf{完备}理性化（Thm 1.5），SARP 再把该理性化的
\textbf{无环性}逼出来（Thm 1.6）。反例 \ref{ex:WARPnotAR} 表明 WARP 只能走到
“完备但可能循环”的理性化——环恰恰藏在 $R^*$ 的不可比对里，只有 SARP 的
链式记账能发现它。这也解释了 Figure 1.1 中 “SARP $\to$ 无环完备理性”
那条不可省略的边。
\end{remark}

\begin{figure}[H]
\centering
\begin{tikzpicture}[
  font=\small,
  box/.style={draw, inner sep=4pt, rounded corners=2pt},
  every node/.style={align=center},
  arr/.style={-{Stealth[length=2.2mm]}, thick}
]
% 上排
\node[box] (harp) at (0,3) {(HARP)};
\node[box] (sarp) at (3.2,3) {(SARP)};
\node[box] (warp) at (6.4,3) {(WARP)};
\node[box] (sca)  at (0,1.2) {(SCA)};
\node[box] (wca)  at (6.4,1.2) {(WCA)};
\draw[arr] (warp) -- (sarp);
\draw[arr] (sarp) -- (harp);
\draw[arr] (wca) -- (sca);
\draw[arr] (harp) -- (sca);
\draw[arr] (sarp) -- node[right=1pt, font=\footnotesize]{无环完备理性} (sca);
\draw[arr] (warp) -- node[right=1pt, font=\footnotesize]{$(R^*\subseteq R_*)$} (wca);
% 下排
\node[box] (fr) at (-3.2,3) {完全理性 full rational};
\node[box] (cr) at (0,-0.4) {完备理性 complete rational};
\node[box] (mr) at (-3.2,-0.4) {M-理性};
\node[box] (mn) at (3.2,-0.4) {(M-正规)};
\node[box] (gn) at (0,-1.8) {(G-正规)};
\node[box] (gr) at (3.2,-1.8) {G-理性};
\draw[arr] (harp) -- (fr);
\draw[arr] (sca) -- (fr);
\draw[arr] (sca) -- (cr);
\draw[arr] (cr) -- (mr);
\draw[arr] (wca) -- (mn);
\draw[arr] (mn) -- (cr);
\draw[arr] (mn) -- (gn);
\draw[arr] (gn) -- (gr);
\draw[arr] (gr) -- (gn);
\draw[arr] (mr) -- (gr);
\end{tikzpicture}
\caption{逻辑蕴含图（重绘自 Suzumura 书 Figure 1.1）。单箭头为不可反转的逻辑蕴含；G-正规 $\Leftrightarrow$ G-理性（Richter 定理 $\gamma$）；HARP $\Leftrightarrow$ SCA $\Leftrightarrow$ 完全理性（Richter--Hansson 定理 $\beta$）。}
\label{fig:logical}
\end{figure}
\figexplan{单箭头为不可反转的逻辑蕴含。上半部分给出显示偏好公理的强弱链：WARP $\Rightarrow$ SARP $\Rightarrow$ HARP，以及 WCA $\Rightarrow$ SCA；下半部分把公理连接到理性化概念（完全理性、完备理性、M-理性、G-理性）。其中 G-正规与 G-理性之间为双向箭头（等价），HARP、SCA 与完全理性三者等价。}

\subsection{反例集：反向箭头为何失败}

以下反例均取 $X=\{x,y,z\}$，是本节必须动手验证的“练习题库”。

\begin{example}[WCA 成立而 SCA 失败]\label{ex:WCAnotSCA}
$\mathcal{S}=\{S_1,S_2,S_3\}$，$S_1=\{x,y\}$，$S_2=\{y,z\}$，$S_3=\{x,z\}$；
$C(S_1)=S_1$，$C(S_2)=S_2$，$C(S_3)=\{x\}$。
显示偏好关系
\[
R^*=\{(x,x),(y,y),(z,z),(x,y),(y,x),(y,z),(z,y),(x,z)\}.
\]
SCA 不满足：$(z,x)\in T(R^*)$（经 $z\mathbin{R^*}y\mathbin{R^*}x$）、$x\in C(S_3)$、
$z\in S_3$，但 $z\notin C(S_3)$。而 WCA 可直接验证成立。
故 SCA $\Rightarrow$ WCA 不可逆。
\end{example}

\begin{example}[G-正规而非 M-正规]\label{ex:GnotM}
$\mathcal{S}=\{S_1,S_2\}$，$S_1=X$，$S_2=\{x,z\}$；$C(S_1)=\{z\}$，$C(S_2)=\{x,z\}$。
此时 $R^*=\{(z,z),(z,x),(z,y),(x,z),(x,x)\}$，
$G^*(S_1)=\{z\}=C(S_1)$、$G^*(S_2)=\{x,z\}=C(S_2)$，故 $C$ 为 G-正规；
但 $M^*(S_1)=\{x,z\}\neq C(S_1)$，故非 M-正规。
\end{example}

\begin{example}[M-理性而不满足 WARP]\label{ex:MnotWARP}
$\mathcal{S}=\{S_1,S_2\}$，$S_1=X$，$S_2=\{y,z\}$；$C(S_1)=\{x,y\}$，$C(S_2)=S_2$。
$R^*=\{(x,x),(y,y),(z,z),(x,y),(x,z),(y,x),(y,z),(z,y)\}$，
可验证 $C$ 为 M-正规（故 M-理性），但
$S_1\cap C(S_2)=\{y,z\}$ 而 $C(S_1)\cap S_2=\{y\}$，WARP 不成立。
\end{example}

\begin{example}[无环完备理性而 SARP 失败]\label{ex:ARnotSARP}
$\mathcal{S}=\{S_1,S_2,S_3\}$，$S_1=\{x,y\}$，$S_2=\{y,z\}$，$S_3=\{x,z\}$；
$C(S_1)=\{x\}$，$C(S_2)=\{y\}$，$C(S_3)=S_3$。
取 $R=\{(x,x),(y,y),(z,z),(x,y),(y,z),(z,x),(x,z)\}$，$R$ 无环且完备，
且 $C(S)=G(S,R)$ 对所有 $S$ 成立——$C$ 是无环完备理性的。
但 $S_1\cap C(S_2)=\{y\}$、$S_2\cap C(S_3)=\{z\}$、$S_3\cap C(S_1)=\{x\}$、
$C(S_1)\cap S_2=\varnothing$、$C(S_2)\cap S_3=\varnothing$，SARP 不满足。
\end{example}

\begin{example}[WARP 不足以保证无环完备理性]\label{ex:WARPnotAR}
同上的集合系，改取 $C(S_1)=\{x\}$，$C(S_2)=\{y\}$，$C(S_3)=\{z\}$。
此时 $R^*=\{(x,y),(y,z),(z,x)\}$，可验证 $R^*\subseteq R_*$（如 $(x,y)\in R_*$
因 $x\in C(S_1)$、$x\notin S_2$ 且 $y\notin C(S_3)$ 等），由定理 \ref{thm:equiv}
的证明知 WARP 成立；但其唯一的完备理性化是循环关系 $x\succ y\succ z\succ x$——
故 WARP 管不住“循环”。
\end{example}

\subsection{理性等级与 Suzumura 一致性闭包}

在任意定义域 $\mathcal{K}_A$ 上，选择函数 $C$ 的\textbf{理性等级}由其理性化关系的强弱定义：

\begin{itemize}
  \item \textbf{理性（R）}：存在理性化 $R$（不加限制）；
  \item \textbf{无环理性（AR）}：存在\textbf{无环}理性化；
  \item \textbf{拟传递理性（QTR）}：存在\textbf{拟传递}理性化；
  \item \textbf{Suzumura 一致理性（SCR）}：存在 \textbf{Suzumura 一致}理性化；
  \item \textbf{完全理性（FR）}：存在是\textbf{序}（完备+传递）的理性化。
\end{itemize}

逻辑关系：
\[
\text{(i)}\ (\mathrm{FR})\Rightarrow(\mathrm{QTR})\Rightarrow(\mathrm{AR})\Rightarrow(\mathrm{R});\qquad
\text{(ii)}\ (\mathrm{FR})\Rightarrow(\mathrm{SCR})\Rightarrow(\mathrm{AR})\Rightarrow(\mathrm{R});
\]
(iii) QTR 与 SCR \textbf{逻辑独立}（反向箭头均不成立）。

Richter--Hansson--Suzumura 的经典理论刻画了两端：FR（HARP/SCA）与 R。
中间的 SCR、QTR、AR 长期悬而未决，直到 Bossert, Sprumont \& Suzumura (2005)
与 Bossert \& Suzumura (2009; 2010) 才补上。其关键工具：

\begin{definition}[Suzumura 一致性闭包]
类比 $T(R)$ 是包含 $R$ 的最小传递超集，定义 $R$ 的 \textbf{Suzumura 一致性闭包}
$SC(R)$ 为包含 $R$ 的最小 Suzumura 一致超集。它有显式表示：
\begin{equation}
SC(R)=R\cup\{(x,y)\in X\times X \mid (x,y)\in T(R)\ \&\ (y,x)\in R\}. \label{eq:SCR}
\end{equation}
\end{definition}

\begin{theorem}[Bossert--Sprumont--Suzumura 2005；Bossert--Suzumura 2010]
任意定义域 $\mathcal{K}_A$ 上的选择函数 $C$ 是 Suzumura 一致理性的（SCR），
当且仅当 $C$ 满足 \textbf{Suzumura 一致闭包一致性公理（SCCC）}——
即把 DRC/TCC 中的显示偏好关系换成其 Suzumura 一致性闭包版本。
\end{theorem}

\noindent SCCC 与 DRC、TCC（Richter 对 R 与 FR 的刻画）在结构上完全平行：
\textbf{每个理性等级对应一个以相应闭包写成的“一致性公理”}——这是显示偏好理论最终完成的形态。

% ============================================================
\section{文献综述：显示偏好理论的研究脉络}
% ============================================================

本节不再以“名录”方式罗列文献，而是把八十年研究脉络\textbf{当作正文来讲}：
每篇经典文献回答三个问题——\textbf{它面对什么缺口、提出了什么公式、留下了什么新问题}。
全部脉络由两条线索交织而成，请始终把握：
\begin{itemize}
  \item \textbf{域的弱化}：预算域 $\mathcal{K}_{CB}$（1938--1950）$\to$
  有限集全域（1959--1973）$\to$ 任意定义域 $\mathcal{K}_A$（1966--2010）；
  \item \textbf{一致性的弱化}：传递 $\to$ 拟传递 $\to$ Suzumura 一致 $\to$ 无环
  （见图 \ref{fig:coherence}）。
\end{itemize}
每弱化一次，旧的刻画定理就失效一次，新的公理（WARP $\to$ SARP $\to$ HARP
$\to$ WCA/SCA $\to$ SCCC）就被发明出来补上缺口——本节即按此逻辑展开。

\begin{figure}[H]
\centering
\begin{tikzpicture}[font=\scriptsize, x=1.12cm]
% 时间轴：1938 -> 0，2010 -> 12（1 单位 = 6 年）
\draw[-{Stealth[length=2.5mm]}, thick] (-0.3,0) -- (12.8,0);
\foreach \x/\yr in {0/1938, 2/1950, 4/1962, 6/1974, 8/1986, 10/1998, 12/2010}
  \draw (\x,0.08) -- (\x,-0.08) node[below=2pt]{\yr};
% ---------- 上方：主线文献（两高度层交错，避免相邻标签相撞）----------
\node[above=13pt, align=center] at (0,0) {Samuelson\\ \scriptsize 1938 WARP};
\draw[gray] (0,0) -- (0,0.28);
\node[above=27pt, align=center] at (2,0) {Houthakker\\ \scriptsize 1950 半传递性};
\draw[gray] (2,0) -- (2,0.62);
\node[above=13pt, align=center] at (3.5,0) {Arrow\\ \scriptsize 1959 抽象域};
\draw[gray] (3.5,0) -- (3.5,0.28);
\node[above=27pt, align=center] at (4.7,0) {Richter\\ \scriptsize 1966 任意域};
\draw[gray] (4.7,0) -- (4.7,0.62);
\node[above=13pt, align=center] at (5.9,0) {Plott\\ \scriptsize 1973 路径无关};
\draw[gray] (5.9,0) -- (5.9,0.28);
\node[above=27pt, align=center] at (7.5,0) {Suzumura\\ \scriptsize 1983 集大成};
\draw[gray] (7.5,0) -- (7.5,0.62);
\node[above=13pt, align=center] at (11.1,0) {Bossert--Sprumont\\ \scriptsize --Suzumura\\ \scriptsize 2005 一致闭包};
\draw[gray] (11.1,0) -- (11.1,0.28);
% ---------- 下方：并行支线（两高度层交错）----------
\node[below=26pt, align=center] at (3.0,0) {Luce 1956\\ \scriptsize 半序（JND）};
\draw[gray] (3.0,0) -- (3.0,-0.60);
\node[below=44pt, align=center] at (5.2,0) {Sen 1969/71\\ \scriptsize 拟传递、$\alpha\beta\gamma$};
\draw[gray] (5.2,0) -- (5.2,-1.05);
\node[below=26pt, align=center] at (6.6,0) {Suzumura 1976/77\\ \scriptsize 一致性、HARP};
\draw[gray] (6.6,0) -- (6.6,-0.60);
\node[below=44pt, align=center] at (9.0,0) {Afriat 1967; Varian 1982\\ \scriptsize 有限数据检验};
\draw[gray] (9.0,0) -- (9.0,-1.05);
\node[below=26pt, align=center] at (11.9,0) {Bossert--Suzumura 2010\\ \scriptsize 等级完成};
\draw[gray] (11.9,0) -- (11.9,-0.60);
% ---------- 阶段带 ----------
\draw[decorate, decoration={brace, amplitude=4pt}, thick] (-0.2,-1.75) -- (2.2,-1.75)
  node[midway, below=5pt, font=\scriptsize]{第一阶段：预算域奠基};
\draw[decorate, decoration={brace, amplitude=4pt}, thick] (10.8,-1.75) -- (12.3,-1.75)
  node[midway, below=5pt, font=\scriptsize]{第四阶段};
\draw[decorate, decoration={brace, amplitude=4pt}, thick] (3.3,-2.40) -- (6.7,-2.40)
  node[midway, below=5pt, font=\scriptsize]{第二阶段：抽象选择域};
\draw[decorate, decoration={brace, amplitude=4pt}, thick] (4.5,-3.00) -- (7.7,-3.00)
  node[midway, below=5pt, font=\scriptsize]{第三阶段：任意域一般理论（与第二阶段交叠）};
\end{tikzpicture}
\caption{显示偏好理论八十年脉络（1938--2010）。注意第二阶段与第三阶段在时间上交叠：
Richter (1966) 的任意域理论早于 Sen (1971) 的有限域工作成熟——两条线索是平行推进、
相互借力的。}
\label{fig:timeline}
\end{figure}
\figexplan{横轴为年份（1938--2010）。上方节点为主线文献：Samuelson 1938 提出 WARP，Houthakker 1950、Arrow 1959、Richter 1966 逐步把理论推向任意域，Suzumura 1976/77 与 1983 集大成；下方节点为并行支线：Luce 1956 半序、Sen 1969/71 拟传递、Afriat 1967 与 Varian 1982 的有限数据检验、Bossert 等 2005/2010。下方花括号标出四个阶段，其中第二阶段与第三阶段在时间上交叠。}

\subsection{第一阶段：预算域上的奠基（1938--1950）}

\textbf{问题背景。}1930 年代的消费者理论建立在可测度效用之上，
但“效用是多少”原则上不可观察。Samuelson 的纲领是：\textbf{只凭可观察的
价格—数量数据，重建消费理论的全部可检验内容}（后来被他称为操作主义
operationalism 在经济学中的落实）。

\textbf{Samuelson (1938, \emph{Economica})} 在竞争预算域
$\mathcal{K}_{CB}=\{B(\boldsymbol{p},M)\}$ 上对单值需求函数
$x^k=c(\boldsymbol{p}^k,M^k)$ 提出两步反反转条件（后称 \textbf{WARP}）：
\begin{equation}
\boldsymbol{p}^0\cdot x^1\le \boldsymbol{p}^0\cdot x^0\ \text{且}\ x^1\neq x^0
\ \Longrightarrow\ \boldsymbol{p}^1\cdot x^0>\boldsymbol{p}^1\cdot x^1.
\label{eq:warpbudget}
\end{equation}
读法：$x^1$ 在第 $0$ 期\textbf{买得起}却没被买（$x^0$ 显示不弱于 $x^1$），
那么当第 $1$ 期实际买了 $x^1$ 时，$x^0$ 必须\textbf{买不起}——否则
$x^1$ 又显示不弱于 $x^0$，两次观察互相打脸。图 \ref{fig:warpviolation}
画出被禁止的构型。

\begin{figure}[H]
\centering
\begin{tikzpicture}[scale=0.95, font=\small]
\draw[-{Stealth[length=2.2mm]}, thick] (0,0) -- (8.2,0) node[below]{$x_1$};
\draw[-{Stealth[length=2.2mm]}, thick] (0,0) -- (0,5.2) node[left]{$x_2$};
% 预算线 B0: (0,4.5)-(5,0); B1: (0,2.6)-(7.5,0)
\draw[very thick] (0,4.5) -- (5,0) node[pos=0.02, above right=-2pt, font=\footnotesize]{$B^0:\ \boldsymbol{p}^0\cdot x=M^0$};
\draw[very thick, dashed] (0,2.6) -- (7.5,0) node[pos=0.55, below right=-2pt, font=\footnotesize]{$B^1:\ \boldsymbol{p}^1\cdot x=M^1$};
\fill (4.2,0.72) circle (2pt) node[below right=1pt]{$x^0=c(B^0)$};
\fill (2.0,1.907) circle (2pt) node[above left=1pt]{$x^1=c(B^1)$};
\fill[gray] (3.434,1.409) circle (1.5pt) node[right=2pt, font=\footnotesize]{交点};
% 标注
\node[font=\footnotesize, align=left] at (6.7,4.0) {$\boldsymbol{p}^0\cdot x^1<\boldsymbol{p}^0\cdot x^0$\\
（第 0 期 $x^1$ 买得起）};
\draw[-{Stealth[length=1.8mm]}, thin] (6.1,3.7) -- (2.35,2.15);
\node[font=\footnotesize, align=left] at (6.7,1.5) {$\boldsymbol{p}^1\cdot x^0<\boldsymbol{p}^1\cdot x^1$\\
（第 1 期 $x^0$ 买得起）};
\draw[-{Stealth[length=1.8mm]}, thin] (6.1,1.35) -- (4.4,0.85);
\end{tikzpicture}
\caption{WARP 禁止的构型：$x^1$ 位于预算线 $B^0$ 内侧（第 0 期买得起却选了 $x^0$），
同时 $x^0$ 位于预算线 $B^1$ 内侧（第 1 期买得起却选了 $x^1$）。
两次观察互相反转，任何完备传递偏好都无法同时解释——这正是 \eqref{eq:warpbudget}
式要排除的情形。}
\label{fig:warpviolation}
\end{figure}
\figexplan{实线为第 0 期预算线 $B^0$，虚线为第 1 期预算线 $B^1$，两线相交。$x^0=c(B^0)$ 与 $x^1=c(B^1)$ 分别是两期的实际选择。关键在于：$x^1$ 落在 $B^0$ 内侧（第 0 期买得起却未选），同时 $x^0$ 落在 $B^1$ 内侧（第 1 期买得起却未选）——两次观察互相反转，任何完备传递偏好都无法同时解释，这正是 WARP 要排除的构型。}

\textbf{留下的缺口}：WARP 在两种商品时足够（Rose 1958, \emph{RES} 证明
$n=2$ 时 WARP $\iff$ 可效用理性化），但 $n\ge 3$ 时不足——
Gale (1960, \emph{Econometrica}) 构造了满足 WARP 却不存在理性化效用的
三商品需求函数。症结在于 WARP 只禁止\textbf{两步}反转，管不住
$x^0\succeq^D x^1\succeq^D x^2\succeq^D x^0$ 式的\textbf{链式}反转
（与本章例 \ref{ex:workedRC} 的“两步绕行”完全同构）。

\textbf{Houthakker (1950, \emph{Economica})} 补上了这个缺口。
记\textbf{直接显示偏好} $x^i\succeq^D x^j \iff
\boldsymbol{p}^i\cdot x^i\ge \boldsymbol{p}^i\cdot x^j$（$x^j$ 在第 $i$ 期买得起），
半传递性公理（semitransitivity，注 \ref{rem:houthakker}）要求：
\begin{equation}
x^0\succeq^D x^1\succeq^D\cdots\succeq^D x^K
\ \Longrightarrow\
\boldsymbol{p}^K\cdot x^K<\boldsymbol{p}^K\cdot x^0
\quad(\text{链上各期观察互不相同}),
\label{eq:semitrans}
\end{equation}
即沿显示偏好链走到 $x^K$ 后，链首 $x^0$ 必须已被\textbf{排除出预算}——
链式反转被一次性禁止。这是预算域上的 SARP。

\textbf{定理（Houthakker 1950）}：需求函数满足半传递性
$\iff$ 存在效用函数 $u$ 使 $c(\boldsymbol{p},M)=
\arg\max\{u(x)\mid \boldsymbol{p}\cdot x\le M\}$。
证明分三步，构成本章 §2.1--§2.5 全部工具的首次合流：
\begin{enumerate}
  \item 由数据定义显示偏好关系 $\succeq^D$ 并取其传递闭包（半传递性保证
  $T(\succeq^D)$ 无反转，即 $T(\succeq^D)$ 的严格部分无环）；
  \item 用 \textbf{Szpilrajn 序扩张}（定理 \ref{thm:szpilrajn}）把该无环关系
  补全为完备序 $R^*$——这是“偏好不可观察”的解决方式：不猜偏好，
  只补全数据给出的部分序；
  \item 对 $R^*$ 调用\textbf{效用表示定理}（Debreu 型，见 §2.2）得 $u$，
  再验证 $\arg\max u$ 与需求一致（生成性，同 §2.5 第三步）。
\end{enumerate}
Houthakker 回忆 Samuelson 收到论文时说：“他不仅补上了我缺失的证明，
而且独立提出了正是我寄望于挽局的那个强公理。”——
显示偏好纲领的\textbf{第一个完备刻画}就此完成。

\subsection{第二阶段：抽象选择域与条件 \texorpdfstring{$\alpha/\beta/\gamma$}{α/β/γ}（1959--1973）}

\textbf{问题转换}。预算集只是“选择面对的集合”的一种特例；
投票、择偶、择校中根本没有价格。能否把显示偏好理论建立在
\textbf{纯集合论}的选择结构上？

\textbf{Arrow (1959, \emph{Econometrica})} 把定义域取为 $X$ 的
\textbf{所有有限子集}（含全部二元、三元组），提出 Arrow 选择公理：
\begin{equation}
S\subseteq T\ \text{且}\ C(T)\cap S\neq\varnothing
\ \Longrightarrow\
C(S)=C(T)\cap S.
\label{eq:arrowaxiom}
\end{equation}
读法：大集合 $T$ 的最优项若在小集合 $S$ 中还有“幸存者”，
则小集合的选择恰好是这些幸存者——不多也不少。
\textbf{Arrow 定理}：$C$ 满足 \eqref{eq:arrowaxiom} $\iff$
存在序 $R$ 使 $C(S)=G(S,R)=\{x\in S\mid (x,y)\in R\ \forall y\in S\}$
对所有 $S$ 成立。这把显示偏好从消费者理论解放为\textbf{一般选择理论}；
代价是域假设极强——现实观察几乎不可能覆盖所有二元组。

\textbf{Sen (1971, \emph{Econometrica})} 把 Arrow 公理拆成
$\alpha$（收缩一致）与 $\beta$（扩张一致）两条（定义 \ref{def:senabc}），
并补上第三条 $\gamma$（并集一致），证明有限集全域上：
\begin{align}
\alpha+\gamma &\iff \exists\ \textbf{完备}\ R:\ C(S)=G(S,R)
\quad(\text{G-理性化}),\\
\alpha+\beta &\iff \exists\ \textbf{序}\ R:\ C(S)=G(S,R)
\quad(\text{完全理性}).
\end{align}
注意 $\alpha$ 在两组刻画中都出现——它是最基本的“菜单无关性”；
$\beta$ 与 $\gamma$ 的分工则精确对应“要不要传递性”：
$\beta$ 控制并列最优者的同进退（这逼出传递性），
$\gamma$ 只要求“常胜者在总决赛仍胜”（这只要完备性）。

\textbf{Sen (1969, \emph{RES})} 同时开辟了弱化传递性的路线。
动机是 Luce 咖啡例 \ref{ex:luce}：不可分辨差异使 $I(R)$ 天然不传递，
而严格偏好 $P(R)$ 的传递性仍可辩护——故只要求\textbf{拟传递}。
他证明：拟传递理性化下社会决策函数（social decision function）
的存在性条件可以系统性给出，集体选择不必卡在 Arrow 不可能性上。

\textbf{Plott (1973, \emph{Econometrica})} 从程序角度补刀传递性。
称选择函数\textbf{路径无关}，若
\begin{equation}
C(S_1\cup S_2)=C\bigl(C(S_1)\cup C(S_2)\bigr),
\label{eq:pathindep}
\end{equation}
即“分批海选再决赛”与“一次性选择”结果相同（这对委员会程序、
淘汰赛设计至关重要）。Plott 证明：路径无关的选择函数可被
\textbf{拟传递}关系刻画，但一般\textbf{不能}被传递关系理性化——
传递性不仅是过强的行为假设，连看似无害的程序合理性都要求放弃它。

\textbf{Blair, Bordes, Kelly \& Suzumura (1976, \emph{JET})}
把上述弱化引入不可能性文献：即使把集体理性从传递放松到
拟传递甚至无环，Arrow 型不可能性定理\textbf{依然成立}——
“放松一致性救不了社会选择”，这反向说明了无环性作为
最后防线的重要性（Bossert--Suzumura 纲领的先声）。

\subsection{第三阶段：任意域上的一般理论（1966--1983）}

\textbf{问题再转换}。Arrow--Sen 要求域含所有有限子集，
这对经验工作仍是奢望：问卷只有几个选项、市场只有几期数据。
能否对定义域\textbf{不作任何限制}？

\textbf{Richter (1966, \emph{Econometrica})} 首次做到这一点。
他直接用显示偏好关系 $R_C$（定义 \ref{def:RC}）“记账”，提出同余公理：
\begin{align}
\text{（D-同余）}\quad &x\in S,\ y\in C(S),\ (x,y)\in R_C
\ \Rightarrow\ x\in C(S);\\
\text{（S-同余）}\quad &x\in S,\ y\in C(S),\ (x,y)\in T(R_C)
\ \Rightarrow\ x\in C(S),
\end{align}
即本章的 WCA 与 SCA（定义 \ref{def:congruence}）。
\textbf{Richter 定理}：任意域上 S-同余 $\iff$ 完全理性；
Richter (1971, \emph{REStud}) 又区分 G-理性化（不要求完备，
用 DRC/TCC 刻画）与 M-理性化，理性化概念的层级由此展开
（§2.6.2 的定理 $\alpha/\beta/\gamma$ 三分）。

\textbf{Hansson (1968, \emph{Theoria})} 独立建立选择结构公理化，
其 C-连通表述（§2.4.2）把“反转”翻译成集合交相等式，
为反例构造提供了最方便的语言——本章 §2.6.3 的反例集全部用此语言书写。

\textbf{Suzumura (1976a, \emph{JET})} 提出介于传递与无环之间的
\textbf{Suzumura 一致性}（定义 \ref{def:weakcoherence}(k)）：
\begin{equation}
(x,y)\in T(R)\ \Rightarrow\ (y,x)\notin P(R).
\end{equation}
其深刻之处：它是使“\textbf{扩张程序}可行”的最弱一致性——
可以证明，$R$ 存在 Suzumura 一致的序扩张 $\iff$ $R$ 本身
Suzumura 一致（这正是 \eqref{eq:SCR} 式闭包 $SC(R)$ 存在的理由）。
传递性在此意义上“强得没必要”，无环性又“弱得不够用”。

\textbf{Suzumura (1977, \emph{Econometrica})} 完成任意域上的主定理：
HARP $\iff$ SCA $\iff$ 完全理性（本章 §2.5 枢纽定理的一般形式）；
(1983) 专著把全部蕴含关系绘成 Figure 1.1（本章图 \ref{fig:logical}），
附录 1.A 给出 Hansson 语言与 Richter 语言的逐步互译。
至此，任意域上“完全理性”一极被彻底封闭。

\subsection{第四阶段：理性等级的完成（2005--2010）}

第三阶段刻画了两端：FR（HARP/SCA）与 R（DRC）。
中间的 AR（无环）、QTR（拟传递）、SCR（Suzumura 一致）
\textbf{悬置了三十年}。困难是技术性的：每个等级的刻画都需要
一个相应的\textbf{闭包算子}，而 $T(R)$ 的构造（定义 \ref{def:TC}）
只对传递有效；拟传递与 Suzumura 一致的“最小超集”是否良定义、
如何显式计算，长期无解。

\textbf{Bossert, Sprumont \& Suzumura (2005, \emph{JET})} 突破第一步，
给出 Suzumura 一致性闭包的\textbf{显式公式}（本章 \eqref{eq:SCR} 式）：
\begin{equation}
SC(R)=R\cup\{(x,y)\in X\times X\mid (x,y)\in T(R)\ \&\ (y,x)\in R\},
\end{equation}
并定义相应同余型公理 \textbf{SCCC}：对任意 $S\in\mathcal{K}$、$x\in S$，
\begin{equation}
\bigl[\forall y\in S:\ (x,y)\in SC(R_C)\bigr]\ \Rightarrow\ x\in C(S).
\end{equation}
\textbf{定理}：任意域上 SCCC $\iff$ SCR。公式 \eqref{eq:SCR} 的直观：
若 $x$ 经链间接不弱于 $y$ 且 $y$ 直接不弱于 $x$，则 $x,y$ 应被
“焊”进同一个无差异类——$SC(R)$ 正是把所有这样的对补上的最小关系。

\textbf{拟传递闭包}（Bossert--Sprumont--Suzumura 2006 前后）：
当 $P(R)$ 无环时，标准构造为
\begin{equation}
(x,y)\in Q(R)\ \Longleftrightarrow\ (y,x)\notin T\bigl(P(R)\bigr),
\qquad P\bigl(Q(R)\bigr)=T\bigl(P(R)\bigr),
\end{equation}
即“严格部分取传递闭包，弱部分取其补”。可验证 $Q(R)$ 拟传递且是
包含 $R$ 的最小拟传递关系；相应公理（QCC）刻画 QTR。
\textbf{无环理性}的刻画最曲折：严格部分的无环扩张必须与选择数据
逐集合对接，Bossert \& Suzumura (2010) 通过扩张程序而非单一闭包公式
给出充要条件。

\textbf{Bossert \& Suzumura (2010)}\ 《Consistency, Choice, and
Rationality》（Harvard University Press）总结整个纲领：
\begin{center}
\emph{每个理性等级 $=$ 一个闭包算子 $+$ 一条以该闭包写成的同余公理。}
\end{center}
至此，图 \ref{fig:logical} 中每个方框都有了\textbf{纯数据可检验}的
充要条件——任意域上理性化的分类学彻底完成。

\subsection{数学基础、教材与应用前沿}

\textbf{序扩张与表示的数学基础。}
Szpilrajn (1930, \emph{Fund.\ Math.}) 的序扩张定理（定理 \ref{thm:szpilrajn}）
是所有“补全偏好”论证的发动机；
Debreu (1954；1959《Theory of Value》) 证明：
\begin{equation}
R\ \text{完备传递且}\textbf{连续}\ (\{y\mid y\succeq x\},\ \{y\mid x\succeq y\}
\ \text{均为闭集})\ \Longrightarrow\ \exists\ \text{连续}\ u\ \text{表示}\ R,
\end{equation}
技术关键是\textbf{可数序稠密子集}的构造（Cantor 方法：
沿例 \ref{ex:countable} 的枚举插入法先把偏好嵌入 $\mathbb{Q}$，再取极限）；
Luce (1956) 半序 $P \iff u(x)>u(y)+\delta$（例 \ref{ex:luce}）说明
无差异不传递不是病态而是知觉阈值的形式化。
Fishburn (1970)《Utility Theory for Decision Making》与
Bridges \& Mehta (1995)《Representations of Preference Orderings》
是表示理论两部标准专著。

\textbf{有限数据检验：从公理到算法。}
Afriat (1967, \emph{IER}) 把 HARP 落到\textbf{有限观察集}
$\{(\boldsymbol{p}^i,x^i)\}_{i=1}^n$：定义
$x^i\succeq^D x^j \iff \boldsymbol{p}^i\cdot x^i\ge\boldsymbol{p}^i\cdot x^j$，
GARP 要求 $T(\succeq^D)$ 无严格反转。
\textbf{Afriat 定理}：GARP $\iff$ 存在\textbf{Afriat 不等式}的解——
效用水平 $U^i$ 与 $\lambda^i>0$ 使
\begin{equation}
U^j\ \le\ U^i+\lambda^i\,\boldsymbol{p}^i\cdot(x^j-x^i)
\qquad(\forall\, i,j),
\label{eq:afriat}
\end{equation}
$\iff$ 存在\textbf{凹、连续、单调}效用函数理性化全部数据。
\eqref{eq:afriat} 是线性不等式组，可直接求解；
Varian (1982, \emph{Econometrica}) 指出检验 GARP 只需算传递闭包
（$O(n^3)$），显示偏好从此成为应用微观（劳动供给、家庭集体模型、
实验经济学）的行业标准检验。Chambers \& Echenique (2016)
《Revealed Preference Theory》（Cambridge）是当代最系统的专著综述，
含计量检验与计算复杂度。

\textbf{教材线索。}
MWG (1995) Chapter 1 是研究生标准入口（选择基/偏好基对偶，
对应本章 §2.3--§2.5，但止于 WARP/SARP）；
Kreps (1988; 2012) 第一章以扩张定理为线索，证明风格最具启发性；
Rubinstein (2012) Lecture 1--3 习题极佳（词典序、半序变体）；
Levin（Stanford Econ 203）讲义 §1--2 与本章对应，含数据检验实例；
本课程 notes3 提供了版式与证明详略标准。

\subsection{脉络总评：一张“缺口—回应”对照表}

\begin{center}
\small
\begin{tabular}{p{4.6cm}p{4.2cm}p{4.4cm}}
\toprule
\textbf{暴露的缺口} & \textbf{回应（公理/工具）} & \textbf{代表文献} \\
\midrule
WARP 只禁两步反转，$n\ge3$ 失效 & 链式半传递性（SARP） & Houthakker 1950 \\
预算域太特殊，无价格场景 & 有限集全域 + Arrow 公理、$\alpha/\beta/\gamma$ & Arrow 1959；Sen 1971 \\
传递性过强（JND、程序设计） & 拟传递、路径无关、无环 & Luce 1956；Sen 1969；Plott 1973 \\
域结构假设不现实 & 同余公理 WCA/SCA、HARP & Richter 1966/71；Hansson 1968 \\
一致性应弱到何处为止 & Suzumura 一致性（扩张可行的最弱条件） & Suzumura 1976a \\
中间理性等级无刻画 & 闭包方法 $SC(R)$、$Q(R)$ & BSS 2005/06；BS 2010 \\
公理无法直接用于有限数据 & GARP + Afriat 不等式 \eqref{eq:afriat} & Afriat 1967；Varian 1982 \\
\bottomrule
\end{tabular}
\end{center}

\begin{remark}[一句话总评]
八十年史 = \textbf{域假设}与\textbf{一致性}的双向弱化史。
每弱化一次，旧定理失效，新公理补位；§2.1--§2.6 的每个定义与定理，
都是对这张对照表中某一行的回应。掌握此表，全章符号皆有归宿。
\end{remark}

% ============================================================
\section{本章小结与课堂检查清单}
% ============================================================

\subsection{一图一表总览}

\begin{center}
\small
\begin{tabular}{llll}
\toprule
\textbf{对象} & \textbf{公理 / 条件} & \textbf{结论} & \textbf{归属} \\
\midrule
偏好 $\succeq$（有限 $X$） & 完备 + 传递 & $\iff$ 存在效用表示 & 经典 \\
偏好 $\succeq$（无限 $X$） & 完备 + 传递 & $\not\Rightarrow$ 表示（词典序反例） & Debreu 需连续性 \\
选择函数 $C$（预算域） & SARP / 半传递性 & $\iff$ 效用最大化理性化 & Houthakker 1950 \\
选择函数 $C$（任意域） & WARP $\iff$ WCA & $\iff$ M-正规（$\Rightarrow$ M-理性） & Suzumura \\
选择函数 $C$（任意域） & HARP $\iff$ SCA & $\iff$ 完全理性（FR） & Richter--Hansson \\
选择函数 $C$（任意域） & G-正规 & $\iff$ G-理性 & Richter \\
选择函数 $C$（任意域） & SCCC & $\iff$ Suzumura 一致理性（SCR） & Bossert--Suzumura \\
\bottomrule
\end{tabular}
\end{center}

\subsection{必须能默写的三个证明}
\begin{enumerate}
  \item 有限集效用表示定理的\textbf{数学归纳}证明（顶层剥皮 + 归纳扩张 $u(x)=n+1$，双方向验证）；
  \item 词典序偏好\textbf{不可表示}的证明（不交非退化区间至多可数）；
  \item HARP $\Rightarrow$ 理性化的\textbf{三步证明}（构造 $\succeq_C$；完备性用 $C_R(\{x,y\})$；
  传递性用 $C_R(\{x,y,z\})$ 三情形；生成性用非空值 + HARP）。
\end{enumerate}

\subsection{课堂提问}
\begin{enumerate}
  \item 若 $C_R(\{a,b\})=\{b\}$，对 $C_R(\{a,b,c\})$ 能推断什么？（$a\notin C_R(\{a,b,c\})$，否则反转）
  \item 效用的序数性如何论证“人际效用比较不可能”？Rawls 与“最小可觉差”各作了什么回应？
  \item 为什么说 M-理性严格强于 G-理性？用 $G(S,R)\subseteq M(S,R)$ 与完备性条件说明。
  \item 例 \ref{ex:WCAnotSCA} 与例 \ref{ex:WARPnotAR} 各自“卡”在哪个逻辑环节？
\end{enumerate}

\vspace{0.6cm}
\noindent\rule{0.35\linewidth}{0.5pt}\par
\noindent{\small 本章 2026 年 9 月整理。数学内容依据 Suzumura《Choice, Preferences, and
Procedures: A Rational Choice Theoretic Approach》Chapter 1（含 Theorem I.1--I.6、
Definition 1.1--1.8、Example 1.1--1.7、Appendix 1.A）与课程材料整理。}

\clearpage
\section{拓展：偏好公理的局部检验与理性化算法}

\subsection{从二元比较到有限数据矩阵}
对有限观测集 $X=\{x^1,\ldots,x^n\}$，可把弱偏好写成矩阵 $A$，其中
$A_{ij}=1$ 当且仅当 $x^i\succeq x^j$。完备性要求每个非对角位置至少有一个方向为 1；传递性要求
$A_{ij}=A_{jk}=1\Rightarrow A_{ik}=1$。矩阵表示使理论公理转化为可编程的逻辑检验。

\begin{figure}[H]\centering
\begin{tikzpicture}[x=1.1cm,y=1.1cm,font=\scriptsize]
\draw[thick] (0,0) rectangle (3.3,3.3);
\foreach \x in {1.1,2.2}{\draw (\x,0)--(\x,3.3);}\foreach \y in {1.1,2.2}{\draw (0,\y)--(3.3,\y);}
\node at (.55,2.75) {$1$};\node at (1.65,2.75) {$1$};\node at (2.75,2.75) {$0$};
\node at (.55,1.65) {$0$};\node at (1.65,1.65) {$1$};\node at (2.75,1.65) {$1$};
\node at (.55,.55) {$1$};\node at (1.65,.55) {$0$};\node at (2.75,.55) {$1$};
\node[above] at (.55,3.3) {$x^1$};\node[above] at (1.65,3.3) {$x^2$};\node[above] at (2.75,3.3) {$x^3$};
\node[left] at (0,2.75) {$x^1$};\node[left] at (0,1.65) {$x^2$};\node[left] at (0,.55) {$x^3$};
\node[align=center] at (1.65,-.55) {列：被比较对象\\行：比较对象};
\end{tikzpicture}
\caption{有限偏好的邻接矩阵}
\end{figure}
\figexplan{矩阵的行列分别对应两个备选方案。对角线上的 1 表示自反性；互为 1 的对称位置表示无差异；单向 1 表示严格偏好。矩阵闭包可用于检验传递性。}

\subsection{传递闭包与算法化检验}
给定直接比较矩阵 $A$，定义其传递闭包 $A^+$：若存在一条从 $i$ 到 $j$ 的有向路径，则 $(A^+)_{ij}=1$。
若 $A^+$ 产生新的反向边而原矩阵没有对应边，则数据包含隐含的严格环或不一致。
Floyd--Warshall 算法可在 $O(n^3)$ 时间内计算闭包：
\[
D^{(k)}_{ij}=D^{(k-1)}_{ij}\lor
(D^{(k-1)}_{ik}\land D^{(k-1)}_{kj}).
\]
这说明“传递性”不仅是抽象公理，也是可执行的有限检验程序。

\begin{proposition}[有限集上的理性化]
若有限选择数据满足无严格偏好环，则存在一个完备且传递的偏好关系理性化该数据；进一步存在一个序数效用函数表示该关系。
\end{proposition}
\begin{proof}
将严格关系视为有向无环图，拓扑排序得到线性序。对线性序赋予整数等级 $1,2,\ldots,n$，即得效用表示。若存在无差异对，则赋予同一等级即可。
\end{proof}

\subsection{WARP、SARP与HARP的层级}
WARP 只排除两期直接反向显示偏好；SARP 进一步排除多期有向环；HARP 将价格—支出不等式纳入检验。
三者不是三个互不相关的定义，而是从局部到全局、从方案关系到预算几何的递进：
\[
\text{WARP}\ \subsetneq\ \text{SARP}\ \subsetneq\ \text{HARP的几何约束}.
\]

\begin{figure}[H]\centering
\begin{tikzpicture}[font=\scriptsize]
\node[draw,rectangle,rounded corners,fill=gray!10,minimum width=2.4cm,minimum height=0.8cm] (w) at (0,0) {WARP：两期};
\node[draw,rectangle,rounded corners,fill=gray!20,minimum width=2.4cm,minimum height=0.8cm] (s) at (4,0) {SARP：多期};
\node[draw,rectangle,rounded corners,fill=gray!30,minimum width=2.8cm,minimum height=0.8cm] (h) at (8,0) {HARP：价格几何};
\draw[->,thick] (w)--(s) node[midway,above] {排除环};\draw[->,thick] (s)--(h) node[midway,above] {加入支出};
\node[align=center] at (4,-1.25) {约束越来越强\\但需要的数据也越来越结构化};
\end{tikzpicture}
\caption{显示偏好公理的递进关系}
\end{figure}
\figexplan{箭头表示检验对象从直接反向选择扩展到多期路径，再扩展到价格与支出构成的几何约束；每一步都增加可证伪内容，但也增加数据要求。}

\subsection{严格偏好环与福利解释}
若 $x\succ y\succ z\succ x$，则不存在任何传递效用函数 $u$ 满足三条严格不等式，因为它们相加后得到 $u(x)>u(x)$。
这类反例说明，效用表示并非对任意排序的机械编码，而是把一致性要求写进数字结构。

\subsection{练习与小结}
\begin{enumerate}
\item 给定四个方案的比较矩阵，计算传递闭包并找出所有严格环。
\item 构造满足 WARP 但不满足 SARP 的三期数据。
\item 说明完备性被放松后，为何单一实值效用函数可能不再是适当表示工具。
\end{enumerate}
本节的核心结论是：偏好理论既可用公理语言表达，也可用矩阵、图和算法检验；二者之间的桥梁正是传递闭包与无环性。

\subsection{核心公理的证明与有限检验}

\begin{proposition}[严格偏好的传递性]
若 $\succeq$ 完备且传递，且 $x\succ y$、$y\succ z$，则 $x\succ z$。
\end{proposition}
\begin{proof}
由严格偏好定义，$x\succeq y$、$y\succeq z$，故传递性给出 $x\succeq z$。若反向关系 $z\succeq x$ 成立，则与 $x\succeq y$ 结合得到 $z\succeq y$，这与 $y\succ z$ 所要求的 $z\not\succeq y$ 矛盾。因此 $z\not\succeq x$，从而 $x\succ z$。
\end{proof}

\begin{proposition}[无差异关系的等价性]
在理性偏好下，$\sim$ 是等价关系，因而可把方案集合分解成互不相交的无差异类。
\end{proposition}
\begin{proof}
自反性来自完备性推出的自反性。对称性直接来自定义。若 $x\sim y$、$y\sim z$，则有 $x\succeq y\succeq z$，传递性给出 $x\succeq z$；同理 $z\succeq x$，故 $x\sim z$。
\end{proof}

\begin{example}[三方案矩阵的逐项检验]
设 $X=\{a,b,c\}$，弱偏好矩阵的非对角元素满足 $a\succeq b$、$b\succeq c$、$a\succeq c$，且反向关系均不成立。则严格偏好为 $a\succ b\succ c$，传递闭包没有新增矛盾边，最优集为 $\{a\}$。若再加入 $c\succeq a$，则传递性迫使 $c\succeq b$，并与 $b\succ c$ 冲突，故该矩阵不能被效用函数表示。
\end{example}

\begin{proposition}[WARP 是 SARP 的二期特例]
在只有两个观测期时，SARP 的显示偏好链只可能形成直接的二环，因此 SARP 检验退化为 WARP 检验。
\end{proposition}
\begin{proof}
两期数据中，若存在反向显示偏好，则有 $x^1\succeq^Rx^2$ 与 $x^2\succeq^Rx^1$。这正是 WARP 禁止的构型，也是长度为二的显示偏好环。除该构型外不存在长度更长的不同期链，因此两期时二者等价。
\end{proof}

\subsection{显示偏好检验的算术展开}
设 $p^t\cdot x^t=m^t$。若 $p^t\cdot x^s\le m^t$，则 $x^s$ 在 $t$ 期可负担，因而 $x^t\succeq^Rx^s$。检验流程必须分为三步：先计算每个交叉支出；再画出有向边；最后寻找返回原点的路径。只检查相邻两期可能漏掉三期环。

\begin{example}[三期环的识别]
若交叉可负担关系为 $x^1\succeq^Rx^2$、$x^2\succeq^Rx^3$、$x^3\succeq^Rx^1$，则传递闭包产生 $x^1\succeq^Rx^1$ 的严格返回路径。若其中至少一条关系严格，SARP 失败；若每条均为无差异，则需要进一步检查是否存在严格环。
\end{example}

\begin{remark}
有限数据检验的逻辑顺序不能倒置：先由价格和数量构造显示偏好，再判断公理，最后讨论是否存在效用表示。直接从观察到的选择宣称“偏好传递”，会把待检验结论当作前提。
\end{remark}

\subsection{显示偏好关系的闭包证明}

\begin{proposition}[传递闭包的最小性]
设 $R$ 是有限方案上的直接显示偏好关系，记 $R^+$ 为所有有限有向路径生成的关系。则 $R^+$ 是包含 $R$ 的最小传递关系。
\end{proposition}
\begin{proof}
任意 $(x,y)\in R$ 是长度为一的路径，故 $R\subseteq R^+$。若 $(x,y),(y,z)\in R^+$，分别存在从 $x$ 到 $y$、从 $y$ 到 $z$ 的有限路径，拼接两条路径即得从 $x$ 到 $z$ 的有限路径，因此 $R^+$ 传递。若 $S$ 是包含 $R$ 的任意传递关系，则对路径长度作归纳：长度一的边在 $S$ 中；若长度 $k$ 的路径对应关系在 $S$ 中，再接一条 $R$ 边，由传递性仍在 $S$ 中。故 $R^+\subseteq S$。
\end{proof}

\begin{example}[闭包如何发现隐含矛盾]
若直接关系只有 $a\succeq^Rb$、$b\succeq^Rc$，则闭包增加 $a\succeq^Rc$。若另有严格边 $c\succ a$，闭包便产生严格环。逐项查看原始边无法发现该矛盾，必须先计算闭包。
\end{example}

\subsection{HARP 检验中的价格—支出不等式}
对观测 $t,s$，若 $p^t\cdot x^s\le p^t\cdot x^t$，则 $x^s$ 在 $t$ 期可负担。若同时 $p^s\cdot x^t<p^s\cdot x^s$，则 $x^t$ 在 $s$ 期严格不可负担，形成严格反向显示偏好。HARP 正是要求这类价格—支出不等式不能沿路径闭合。

\begin{proposition}[支出不等式的尺度不变性]
若每期价格与收入同时乘以同一正数 $a_t$，则 GARP/HARP 的可负担关系不变。
\end{proposition}
\begin{proof}
$a_tp^t\cdot x^s\le a_tp^t\cdot x^t$ 当且仅当 $p^t\cdot x^s\le p^t\cdot x^t$，因为 $a_t>0$ 不改变不等式方向。故所有直接显示偏好边及其闭包均不变。
\end{proof}

\begin{remark}
该尺度不变性解释了为什么显示偏好检验只依赖相对价格与相对支出，而不依赖货币单位。把元改为分、把收入和价格同时乘以一百，不会改变理性化结论。
\end{remark}

\subsection{理性化命题的证明细节}

\begin{proposition}[理性化与最大元]
若 $\succeq$ 完备且传递，且 $B$ 有限非空，则
$C(B)=\{x\in B:x\succeq y,\ \forall y\in B\}$ 非空。
\end{proposition}
\begin{proof}
从 $B$ 中任取 $x_1$。若它击败所有元素，则完成；否则存在 $x_2$ 使 $x_2\succ x_1$。继续比较。由于 $B$ 有限，严格改进不可能无限进行，过程在某个 $x_k$ 停止。若仍存在 $y\succ x_k$，应继续改进，矛盾。因此 $x_k$ 是最大元，且属于选择集。
\end{proof}

\begin{example}[并列最大元]
若 $a\sim b\succ c$，则 $C(\{a,b,c\})=\{a,b\}$。理性化并不意味着选择对应必须单值；只有严格偏好或严格拟凹等附加条件才能保证唯一选择。
\end{example}

\begin{proposition}[严格偏好环排除理性化]
若有限选择数据的显示偏好闭包包含严格环，则不存在完备传递偏好关系理性化该数据。
\end{proposition}
\begin{proof}
设严格环为 $x^1\succ x^2\succ\cdots\succ x^k\succ x^1$。严格偏好的传递性推出 $x^1\succ x^1$，违反严格偏好的反自反性。因此任何理性偏好都不可能包含该环。
\end{proof}

\subsection{从选择数据到效用不等式}
若 $x^t$ 在预算 $B^t$ 中被选择，且 $x^s$ 也可负担，则理性化要求 $u(x^t)\ge u(x^s)$。一串可负担关系便产生一串效用弱不等式；若最后回到起点并至少有一次严格不可负担，则得到 $u(x^t)>u(x^t)$ 的矛盾。这是 GARP 的算术核心。

\begin{example}[一条三期不等式链]
若 $x^1$ 在第 1 期可负担 $x^2$，$x^2$ 在第 2 期可负担 $x^3$，而第 3 期严格排除 $x^1$，则理性化要求
$u(x^1)\ge u(x^2)\ge u(x^3)>u(x^1)$，显然不可能。该链条比单独观察任意一对数据更有信息。
\end{example}

\begin{remark}
显示偏好定理的方向必须严格区分：数据满足公理，推出存在某类效用函数；数据违反公理，推出不存在满足相应性质的理性化。不能把“某个简单效用函数拟合失败”误写成“所有效用函数都无法理性化”。
\end{remark}

\subsection{GARP 不等式链的完整展开}
令第 $t$ 期观察到的选择为 $x^t$，价格为 $p^t$，收入为 $m^t=p^t\cdot x^t$。若 $x^s$ 在第 $t$ 期可负担，则
\[
 p^t\cdot x^s\le p^t\cdot x^t
 \quad\Longrightarrow\quad x^t\succeq x^s.
\]
若沿若干期得到 $x^1\succeq x^2\succeq\cdots\succeq x^k$，且最后一期严格显示排除 $x^1$，则理性化要求
\[
 u(x^1)\ge u(x^2)\ge\cdots\ge u(x^k)>u(x^1),
\]
矛盾。这个推导将图形中的“可负担边”转化为效用数值上的不等式链；GARP 的核心不是某一条价格比较，而是所有可负担边的传递闭包。

\begin{example}[三期链的逐项计算]
假设 $p^1\cdot x^2\le p^1\cdot x^1$、$p^2\cdot x^3\le p^2\cdot x^2$，且
$p^3\cdot x^1<p^3\cdot x^3$。前两式分别推出 $u(x^1)\ge u(x^2)$ 与
$u(x^2)\ge u(x^3)$；若数据还含有从 $x^3$ 回到 $x^1$ 的显示弱偏好链，便得到严格闭环。故检查时必须把“直接边”和“闭包边”分开记录，避免把一条不可负担关系误读为偏好关系。
\end{example}

\begin{proposition}[尺度变换不改变显示偏好]
若每期价格与收入同时乘以同一 $a_t>0$，则预算可负担关系及 WARP、SARP、GARP 的判定均不变。
\end{proposition}
\begin{proof}
$(a_tp^t)\cdot x^s\le a_tm^t$ 当且仅当 $p^t\cdot x^s\le m^t$。因此每一条可负担边保持不变，由这些边生成的传递闭包也保持不变，所有环和严格边的存在性自然不变。
\end{proof}

\subsection{有限数据的手算程序与逻辑检查}
对有限选择数据，理性化检验可以按四列进行：第一列记录观察期 $t$；第二列记录被选择束 $x^t$；第三列计算每个其他束在该期的支出 $p^t\cdot x^s$；第四列把可负担关系转化为有向边 $t\to s$。完成表格后，先查两期反向边，再查传递闭包中的环，最后检查环上是否至少存在一条严格边。没有严格边的闭环只说明无差异，不构成 GARP 违反；有严格边的闭环才推出效用不等式矛盾。

设三期数据的支出矩阵为 $E_{ts}=p^t\cdot x^s$。对角线 $E_{tt}=m^t$ 是实际支出。若 $E_{ts}\le E_{tt}$，就在矩阵中标记 $t\to s$；若不等号严格且 $x^s\ne x^t$，再标记为严格边。于是，GARP 的计算对象不是价格矩阵本身，而是由 $E$ 的逐行比较生成的标记矩阵。

\begin{example}[四期闭包的检查顺序]
若直接边为 $1\to2$、$2\to3$、$3\to4$，算法先生成传递边 $1\to3$、$1\to4$、$2\to4$。若另有 $4\to1$，则出现闭环；此时只需检查四条边中是否有一条严格边。若有，便得到 $u(x^1)\ge u(x^2)\ge u(x^3)\ge u(x^4)>u(x^1)$；若没有，四个束可以属于同一无差异类。将“闭包”与“严格性”分开，是避免把普通循环误判为违反的关键。
\end{example}

\subsection{从 WARP 到 SARP 再到 GARP}
WARP 只检查两期直接反转，SARP 检查直接边与传递边形成的严格循环，GARP 则允许无差异并只排除带严格边的闭环。因此三者的差别不在于是否使用价格，而在于所允许的闭包范围和对并列关系的处理方式。两期数据中不存在更长路径，故 SARP 的检验退化为 WARP；三期及以上时，单独检查每一对边不再足够。

若数据满足 GARP，Afriat 不等式给出 $U_t\le U_s+\lambda_s p^s\cdot(x^t-x^s)$，其中 $\lambda_s>0$。逐项解出 $U_t$ 与 $\lambda_s$ 后，可构造分段线性凹效用 $u(x)=\min_s\{U_s+\lambda_s p^s\cdot(x-x^s)\}$。在观测点 $x^t$ 处，第 $t$ 个平面取值为 $U_t$，其余平面不低于 $U_t$，所以 $u(x^t)=U_t$；在预算内任意束都不会超过该期观测效用。

\subsection{三期数据的逐项理性化核对}
设三期观测为 $(p^t,x^t)$，并假定每期选择都耗尽预算。对任意两期 $s,t$，先计算
\[
 p^s\cdot x^t-p^s\cdot x^s.
\]
若该差值不大于零，则 $x^t$ 在第 $s$ 期可负担，观察到的选择意味着
$x^s\succeq x^t$；若差值严格小于零且 $x^t\ne x^s$，则显示偏好为严格关系。
因此，理性化检验不能只看价格向量是否变化，而必须逐期计算“另一期选择在本期预算中的支出”。

例如，若
\[
 p^1=(1,2),\quad x^1=(2,1),\qquad
 p^2=(2,1),\quad x^2=(1,2),
\]
则
\[
 p^1\cdot x^1=4,\quad p^1\cdot x^2=5;
 \qquad
 p^2\cdot x^2=4,\quad p^2\cdot x^1=5.
\]
两期互不显示可负担，故不存在由这两期单独构成的直接环。若第三期选择
$x^3$ 同时落入前两期预算约束，就必须把三条弱不等式连起来检查；一旦得到
$u(x^1)\ge u(x^2)\ge u(x^3)>u(x^1)$，矛盾来自闭环，而不是来自任意一条单独比较。

\begin{proposition}[尺度变换不改变显示偏好]
若每期价格与收入同时乘以同一 $a_t>0$，则预算可负担关系及 WARP、SARP、GARP 的判定均不变。
\end{proposition}
\begin{proof}
$(a_sp^s)\cdot x^t\le a_sm^s$ 当且仅当 $p^s\cdot x^t\le m^s$。因此每一条可负担边都保持不变，
由这些边生成的传递闭包也保持不变，所有环和严格边的存在性自然不变。
\end{proof}


\end{document}
